% Przewodnik seminaryjny.
%
% Dokument sklada sie ta sama klasa, co prace studentow, i korzysta z tego samego
% stylu bibliograficznego. Jest wiec jednoczesnie instrukcja i pelnowymiarowym
% przykladem dzialania wzoru: kazdy element skladu opisany w rozdziale pierwszym
% wystepuje dalej w praktyce.
%
% Kompilacja z katalogu przewodnik/:
%   latexmk -lualatex przewodnik.tex
% Plik latexmkrc wskazuje katalog wzoru, w ktorym lezy klasa i styl bibliograficzny.

\documentclass{epi-praca}

\addbibresource{bibliografia/zrodla.bib}

\autor{dr Marek Deja}
\tytul{Przewodnik seminaryjny}
\podtytul{Praca licencjacka i projekt dyplomowy na kierunku EPI}
\rok{2027}

\begin{document}

\stronatytulowaporadnika

\tableofcontents

\wprowadzenie

Przewodnik obejmuje cały cykl seminarium: od pierwszego uruchomienia wzoru, przez budowę pakietu i pisanie kolejnych rozdziałów, po kontrolę formalną przed oddaniem pracy.

Praca dyplomowa na kierunku elektroniczne przetwarzanie informacji składa się z dwóch ściśle powiązanych części: projektu dyplomowego oraz pracy licencjackiej, która ten projekt opisuje. Zaliczenie seminarium następuje dopiero po przyjęciu projektu \parencite{standardyEPI}. Twoim projektem jest pakiet języka R implementujący algorytm z artykułu metodycznego wraz z serwisem, który go udostępnia.

\section{Dla kogo i po co}

Przewodnik zakłada, że znasz podstawy statystyki i programowania w R \parencite{rcore}, a nie znasz LaTeX-a ani inżynierii oprogramowania. Wszystkiego, co potrzebne, nauczysz się tutaj od zera.

Dokument, który czytasz, jest jednocześnie instrukcją i przykładem. Złożono go tą samą klasą, z której korzystasz przy pisaniu pracy, i tym samym stylem bibliograficznym. Każdy element opisany w rozdziale pierwszym -- tabela z tytułem nad nią, wzór z etykietą, listing kodu, cytowanie w nawiasie, spis ilustracji -- występuje dalej w praktyce. Kiedy nie pamiętasz zapisu, zajrzyj do źródła tego przewodnika: leży w tym samym repozytorium.

\section{Jak czytać}

Rozdziały pierwszy, drugi i trzeci dotyczą narzędzi. Przeczytaj je na początku seminarium, zanim napiszesz pierwsze zdanie pracy. Rozdział pierwszy wystarczy przerobić do połowy, żeby zacząć -- reszta przyda się później.

Rozdziały czwarty, piąty i szósty dotyczą pisania. Do rozdziału czwartego wracaj przy każdej kolejnej części pracy. Rozdział piąty przeczytaj raz, na samym początku, bo pytanie badawcze stawia się przed pisaniem, a nie po nim. Rozdział szósty służy jako lista kontrolna przed oddaniem.

Rozdziały siódmy i ósmy dotyczą pracy z narzędziem wspomagającym pisanie kodu. Siódmy opisuje metodę, ósmy granice dozwolonego użycia i sposób jego dokumentowania. Oba przeczytaj \wyroznienie{zanim wydasz pierwsze polecenie}, a nie wtedy, gdy pakiet będzie już gotowy: część wymagań dotyczy rzeczy, których nie da się odtworzyć po fakcie.

\begin{tabelaepi}{Zawartość przewodnika}{tab:zawartosc}
\begin{tabular}{@{}cp{6.4cm}p{6.2cm}@{}}
\toprule
Rozdział & Temat & Kiedy czytać \\
\midrule
1 & LaTeX & na początku, pierwsze dwa poziomy \\
2 & Bibliografia i Zotero & przed pierwszą lekturą \\
3 & Overleaf i współpraca & przy zakładaniu projektu \\
4 & Jak napisać pracę & przy pisaniu każdej części \\
5 & Warsztat akademicki & raz, na samym początku \\
6 & Wymogi formalne & przed oddaniem \\
7 & Praca z agentem programistycznym & przed pierwszym poleceniem \\
8 & Jawność i odpowiedzialność & na początku i przed oddaniem \\
\bottomrule
\end{tabular}
\end{tabelaepi}

\section{Dokumenty rozstrzygające}

Wymagania wobec pracy określają dwa dokumenty Instytutu. \wyroznienie{Standardy prac dyplomowych EPI} \parencite{standardyEPI} rozstrzygają o układzie pracy i o wymogach wobec projektu dyplomowego. \wyroznienie{Instrukcja ISI} \parencite{instrukcjaISI} rozstrzyga o składzie, przypisach i opisach bibliograficznych. Tam, gdzie dokumenty się rozchodzą, w przewodniku wskazany jest ten, który obowiązuje.

Wymagania te są we wzorze zrealizowane automatycznie. Nie zmieniaj ustawień składu -- rozbieżność z wymogami obciąża pracę, a nie wzór.

\section{Dwie ścieżki pracy}\label{sec:sciezki}

Pracę można pisać na dwa sposoby i przewodnik opisuje oba.

\wyroznienie{Ścieżka przeglądarkowa.} Piszesz na uczelnianej instalacji Overleaf.
Nie instalujesz niczego, kompilacja dzieje się na serwerze, promotor widzi tę samą
wersję co Ty. Jest to ścieżka domyślna i wystarczająca do napisania całej pracy.

\wyroznienie{Ścieżka lokalna.} Piszesz u siebie, w Positronie, mając zainstalowaną
dystrybucję \TeX-a i R. Zyskujesz to, że praca, pakiet i wykresy są w jednym miejscu,
oraz dostęp do skryptów kontrolnych, których na Overleaf uruchomić się nie da.

Różnica jest istotna w kilku miejscach: przy kompilacji, przy czyszczeniu plików
pomocniczych i przy kontroli objętości. Wszędzie tam, gdzie postępowanie się rozchodzi,
akapity są oznaczone słowami \wyroznienie{Lokalnie} i \wyroznienie{Na Overleaf}.
Tam, gdzie takiego oznaczenia nie ma, opis dotyczy obu ścieżek.

Najczęstszy układ w praktyce jest mieszany: pakiet i wykresy powstają lokalnie,
tekst pracy na Overleaf. Wtedy stosujesz oznaczenie właściwe dla czynności, a nie
dla całej pracy.

\chapter{LaTeX od podstaw}\label{ch:latex}

W rozdziale poznajesz LaTeX-a na wzorze, którym piszesz pracę, a nie w oderwaniu od niego. Wszystkie przykłady działają bez dodatkowej konfiguracji, bo ustawienia składu są w \termin{klasie dokumentu}, czyli w pliku, w którym zapisany jest cały wygląd pracy. Twoja klasa to \plik{epi-praca.cls} i nie musisz do niej zaglądać.

Podziel naukę na trzy wieczory. Po pierwszym umiesz złożyć tekst z rozdziałami i cytowaniami. Po drugim -- tabele, rysunki i wzory. Trzeci przyda się dopiero przy redakcji.

\section{Jak myśleć o LaTeX-u}

W edytorze tekstu zaznaczasz fragment i nadajesz mu wygląd. W LaTeX-u \wyroznienie{opisujesz, czym fragment jest}, a wygląd wynika z tego automatycznie. Piszesz \enquote{to jest tytuł rozdziału}, a nie \enquote{to ma być osiemnaście punktów, wytłuszczone, wyśrodkowane}.

Ma to trzy konsekwencje, które w pracy licencjackiej są warte więcej niż koszt nauki. Skład jest zgodny z wymogami Instytutu z definicji, a nie dlatego, że pamiętałeś o ustawieniu interlinii w każdym akapicie. Spis treści, spis ilustracji, bibliografia i indeks nazwisk powstają same z tego, co napisałeś w tekście -- nie da się mieć w spisie treści rozdziału, którego nie ma. Wzory matematyczne są zapisem, a nie obrazkiem, co przy pracy pełnej wzorów z artykułu metodycznego przestaje być wygodą, a staje się warunkiem wykonalności.

Kosztem jest to, że pliku źródłowego się nie ogląda -- trzeba go najpierw \termin{skompilować}, czyli przetworzyć w gotowy dokument PDF. Do tego przyzwyczaisz się w pierwszy dzień.

\section{Poziom pierwszy: tekst i struktura}

\subsection{Anatomia dokumentu}

Praca składa się z pliku głównego i plików rozdziałów.

\begin{zapis}
\documentclass{epi-praca}        % klasa: caly sklad

\addbibresource{bibliografia/literatura.bib}
\autor{Jan Kowalski}             % metadane
\tytul{Tytul pracy}

\begin{document}                 % poczatek tresci
\stronatytulowa
\tableofcontents
\wprowadzenie

Przewodnik obejmuje cały cykl seminarium: od pierwszego uruchomienia wzoru, przez budowę pakietu i pisanie kolejnych rozdziałów, po kontrolę formalną przed oddaniem pracy.

Praca dyplomowa na kierunku elektroniczne przetwarzanie informacji składa się z dwóch ściśle powiązanych części: projektu dyplomowego oraz pracy licencjackiej, która ten projekt opisuje. Zaliczenie seminarium następuje dopiero po przyjęciu projektu \parencite{standardyEPI}. Twoim projektem jest pakiet języka R implementujący algorytm z artykułu metodycznego wraz z serwisem, który go udostępnia.

\section{Dla kogo i po co}

Przewodnik zakłada, że znasz podstawy statystyki i programowania w R \parencite{rcore}, a nie znasz LaTeX-a ani inżynierii oprogramowania. Wszystkiego, co potrzebne, nauczysz się tutaj od zera.

Dokument, który czytasz, jest jednocześnie instrukcją i przykładem. Złożono go tą samą klasą, z której korzystasz przy pisaniu pracy, i tym samym stylem bibliograficznym. Każdy element opisany w rozdziale pierwszym -- tabela z tytułem nad nią, wzór z etykietą, listing kodu, cytowanie w nawiasie, spis ilustracji -- występuje dalej w praktyce. Kiedy nie pamiętasz zapisu, zajrzyj do źródła tego przewodnika: leży w tym samym repozytorium.

\section{Jak czytać}

Rozdziały pierwszy, drugi i trzeci dotyczą narzędzi. Przeczytaj je na początku seminarium, zanim napiszesz pierwsze zdanie pracy. Rozdział pierwszy wystarczy przerobić do połowy, żeby zacząć -- reszta przyda się później.

Rozdziały czwarty, piąty i szósty dotyczą pisania. Do rozdziału czwartego wracaj przy każdej kolejnej części pracy. Rozdział piąty przeczytaj raz, na samym początku, bo pytanie badawcze stawia się przed pisaniem, a nie po nim. Rozdział szósty służy jako lista kontrolna przed oddaniem.

Rozdziały siódmy i ósmy dotyczą pracy z narzędziem wspomagającym pisanie kodu. Siódmy opisuje metodę, ósmy granice dozwolonego użycia i sposób jego dokumentowania. Oba przeczytaj \wyroznienie{zanim wydasz pierwsze polecenie}, a nie wtedy, gdy pakiet będzie już gotowy: część wymagań dotyczy rzeczy, których nie da się odtworzyć po fakcie.

\begin{tabelaepi}{Zawartość przewodnika}{tab:zawartosc}
\begin{tabular}{@{}cp{6.4cm}p{6.2cm}@{}}
\toprule
Rozdział & Temat & Kiedy czytać \\
\midrule
1 & LaTeX & na początku, pierwsze dwa poziomy \\
2 & Bibliografia i Zotero & przed pierwszą lekturą \\
3 & Overleaf i współpraca & przy zakładaniu projektu \\
4 & Jak napisać pracę & przy pisaniu każdej części \\
5 & Warsztat akademicki & raz, na samym początku \\
6 & Wymogi formalne & przed oddaniem \\
7 & Praca z agentem programistycznym & przed pierwszym poleceniem \\
8 & Jawność i odpowiedzialność & na początku i przed oddaniem \\
\bottomrule
\end{tabular}
\end{tabelaepi}

\section{Dokumenty rozstrzygające}

Wymagania wobec pracy określają dwa dokumenty Instytutu. \wyroznienie{Standardy prac dyplomowych EPI} \parencite{standardyEPI} rozstrzygają o układzie pracy i o wymogach wobec projektu dyplomowego. \wyroznienie{Instrukcja ISI} \parencite{instrukcjaISI} rozstrzyga o składzie, przypisach i opisach bibliograficznych. Tam, gdzie dokumenty się rozchodzą, w przewodniku wskazany jest ten, który obowiązuje.

Wymagania te są we wzorze zrealizowane automatycznie. Nie zmieniaj ustawień składu -- rozbieżność z wymogami obciąża pracę, a nie wzór.

\section{Dwie ścieżki pracy}\label{sec:sciezki}

Pracę można pisać na dwa sposoby i przewodnik opisuje oba.

\wyroznienie{Ścieżka przeglądarkowa.} Piszesz na uczelnianej instalacji Overleaf.
Nie instalujesz niczego, kompilacja dzieje się na serwerze, promotor widzi tę samą
wersję co Ty. Jest to ścieżka domyślna i wystarczająca do napisania całej pracy.

\wyroznienie{Ścieżka lokalna.} Piszesz u siebie, w Positronie, mając zainstalowaną
dystrybucję \TeX-a i R. Zyskujesz to, że praca, pakiet i wykresy są w jednym miejscu,
oraz dostęp do skryptów kontrolnych, których na Overleaf uruchomić się nie da.

Różnica jest istotna w kilku miejscach: przy kompilacji, przy czyszczeniu plików
pomocniczych i przy kontroli objętości. Wszędzie tam, gdzie postępowanie się rozchodzi,
akapity są oznaczone słowami \wyroznienie{Lokalnie} i \wyroznienie{Na Overleaf}.
Tam, gdzie takiego oznaczenia nie ma, opis dotyczy obu ścieżek.

Najczęstszy układ w praktyce jest mieszany: pakiet i wykresy powstają lokalnie,
tekst pracy na Overleaf. Wtedy stosujesz oznaczenie właściwe dla czynności, a nie
dla całej pracy.

\end{document}                   % koniec tresci
\end{zapis}

Wszystko przed \kod{\textbackslash begin\{document\}} to \termin{preambuła}, czyli ustawienia. Wszystko między \kod{\textbackslash begin\{document\}} a \kod{\textbackslash end\{document\}} to treść. W Twojej pracy preambuła jest już gotowa; uzupełniasz w niej wyłącznie metadane.

\wyroznienie{Polecenie} zaczyna się od ukośnika wstecznego i mówi, co ma zostać zrobione. To, na czym polecenie działa, nazywa się \termin{argumentem}. Argumenty obowiązkowe są w nawiasach klamrowych, opcjonalne w kwadratowych: w zapisie \kod{\textbackslash includegraphics[width=5cm]\{wykres.png\}} nazwa pliku jest argumentem obowiązkowym, a szerokość opcjonalnym.

\wyroznienie{Środowisko} obejmuje cały fragment tekstu i ma jawny początek oraz koniec. Otwiera je \kod{\textbackslash begin} z nazwą, zamyka \kod{\textbackslash end} z tą samą nazwą, a wszystko pomiędzy składa się według reguł tego środowiska -- jak \kod{quote} dla cytatu blokowego albo \kod{tabular} dla tabeli. Środowiskiem jest też sam dokument, o czym mówi para poleceń z przykładu wyżej.

\subsection{Akapity i spacje}

Nowy akapit zaczyna się po \wyroznienie{pustej linii}. Pojedyncze złamanie wiersza w pliku źródłowym nie znaczy nic -- LaTeX i tak złoży tekst od nowa. Wielokrotne spacje redukują się do jednej. Nie wymuszaj odstępów spacjami ani pustymi liniami; od tego są polecenia, a w Twojej pracy odstępy ustawia klasa.

W seminarium obowiązuje jedna zasada zapisu źródła: \wyroznienie{jeden akapit w jednym wierszu}. Akapitu nie łamiemy w pliku, bez względu na jego długość; wiersze rozdziela wyłącznie pusta linia. Powód jest praktyczny -- porównanie wersji pokazuje wtedy, który akapit się zmienił, zamiast lawiny przesuniętych wierszy po dopisaniu jednego słowa. Ma to znaczenie przy pracy z promotorem i w historii repozytorium, a edytory i tak zawijają długie wiersze na ekranie. Zasada nie dotyczy wnętrza listingów, tabel i wzorów ani pozycji list wypunktowanych.

\wyroznienie{Twarda spacja} zapisywana tyldą nie pozwala złamać wiersza w danym miejscu. Używaj jej tam, gdzie przeniesienie wygląda źle:

\begin{zapis}
rysunek~\ref{rys:mapa}    s.~15-21    prof.~Kowalski    nr~3
\end{zapis}

\subsection{Polskie znaki i typografia}

Piszesz normalnie: zażółć gęślą jaźń. Plik musi być zapisany w UTF-8; Overleaf i Positron robią to domyślnie.

\begin{tabelaepi}{Zapis typograficzny}{tab:typografia}
\begin{tabular}{@{}lll@{}}
\toprule
Chcesz & Piszesz & Efekt \\
\midrule
cudzysłów polski & \kod{\textbackslash enquote\{cytat\}} & \enquote{cytat} \\
cytat w cytacie & \kod{\textbackslash enquote\{a \textbackslash enquote\{b\} c\}} & \enquote{a \enquote{b} c} \\
półpauza & \kod{-{}-} & -- \\
łącznik & \kod{-} & - \\
\bottomrule
\end{tabular}
\end{tabelaepi}

Cudzysłowów nie wpisuj ręcznie. Polecenie \kod{\textbackslash enquote} dobiera je do języka i poprawnie zagnieżdża, dokładnie tak, jak przewiduje Instrukcja \parencite{instrukcjaISI}.

W seminarium obowiązuje jedna zasada typograficzna wykraczająca poza wymogi Instytutu: \wyroznienie{nie stosujemy pauzy}, czyli trzech łączników. Zdanie wtrącone wydzielamy półpauzą z odstępem po obu stronach -- tak jak w tym zdaniu. Łącznik pozostaje łącznikiem w wyrazach złożonych i w zakresach stron.

Kilka znaków ma w LaTeX-u znaczenie specjalne i żeby pojawiły się w tekście, trzeba je poprzedzić ukośnikiem: procent, dolar, ampersand, podkreślenie, kratka i nawiasy klamrowe. Najczęściej potykasz się o podkreślenie w nazwach funkcji -- dlatego do nazw służą polecenia opisane w podrozdziale~\ref{sec:makra}, które robią to za Ciebie.

\subsection{Rozdziały i podrozdziały}

\begin{zapis}
\chapter{Podstawy metodyczne}          % 1
\section{Aparat formalny}              % 1.1
\subsection{Normalizacja}              % 1.1.1
\end{zapis}

Numeracja jest automatyczna. Nie wpisuj numerów ręcznie -- po wstawieniu rozdziału w środku przenumerowanie zajmie sekundę, a nie godzinę. Elementy nienumerowane, obecne mimo to w spisie treści, mają własne polecenia: \kod{\textbackslash wprowadzenie}, \kod{\textbackslash podsumowanie}, \kod{\textbackslash wykazzrodel}, \kod{\textbackslash aneksy}.

\subsection{Etykiety i odsyłacze}

To jest funkcja, dla której warto uczyć się LaTeX-a. Nadajesz obiektowi etykietę i odwołujesz się do niej po nazwie; numer wstawia się sam i zawsze jest poprawny.

\begin{zapis}
\section{Aparat formalny}\label{sec:aparat}

Metode omowiono w podrozdziale~\ref{sec:aparat} na stronie~\pageref{sec:aparat}.
\end{zapis}

\begin{tabelaepi}{Konwencja przedrostków w etykietach}{tab:przedrostki}
\begin{tabular}{@{}ll@{}}
\toprule
Przedrostek & Obiekt \\
\midrule
\kod{sec:} & rozdział, podrozdział \\
\kod{tab:} & tabela \\
\kod{rys:} & rysunek \\
\kod{wyk:} & wykres \\
\kod{lst:} & listing \\
\kod{eq:} & wzór \\
\bottomrule
\end{tabular}
\end{tabelaepi}

Odsyłacz do nieistniejącej etykiety daje w gotowym pliku dwa znaki zapytania i ostrzeżenie \kod{Reference ... undefined}. Zawsze sprawdzaj je przed oddaniem.

\subsection{Cytowania}

\begin{zapis}
\parencite{cisek2002filozoficzne}                  % (Cisek 2002)
\parencite[s.~15-21]{kowalski2010a}                % (Kowalski 2010a, s. 15-21)
\parencite{klucz1,klucz2}                          % (A 2019; B 2020)
\textcite{wozniak1997kognitywizm} pokazuje, ze...  % Wozniak (1997) pokazuje, ze...
\end{zapis}

Postać cytowania -- inicjał przy zbieżnych nazwiskach, skrót \enquote{i in.} przy więcej niż trzech autorach, sufiksy rocznika przy tym samym roku -- dobiera styl automatycznie. Twoim zadaniem jest poprawny wpis w pliku bibliograficznym, opisany w rozdziale~\ref{ch:bibliografia}.

\subsection{Polecenia własne wzoru}\label{sec:makra}

Nie sięgaj po polecenia formatujące bezpośrednio. Wzór ma polecenia, które mówią, \wyroznienie{czym} coś jest, i dobierają wygląd zgodnie z wymogami Instytutu.

\begin{tabelaepi}{Polecenia semantyczne wzoru}{tab:makra}
\begin{tabular}{@{}llp{5.4cm}@{}}
\toprule
Polecenie & Zastosowanie & Wygląd \\
\midrule
\kod{\textbackslash fun\{przygotuj\_dane\}} & nazwa funkcji & \fun{przygotuj\_dane} \\
\kod{\textbackslash argument\{skala\}} & nazwa argumentu & \argument{skala} \\
\kod{\textbackslash pkg\{ggplot2\}} & nazwa pakietu & \pkg{ggplot2} \\
\kod{\textbackslash plik\{R/algorytm.R\}} & plik lub katalog & \plik{R/algorytm.R} \\
\kod{\textbackslash kod\{x <- 1\}} & krótki kod w zdaniu & \kod{x <- 1} \\
\kod{\textbackslash termin\{fuzzyfikacja\}} & termin, wyraz obcy, tytuł & \termin{fuzzyfikacja} \\
\kod{\textbackslash wyroznienie\{ważne\}} & wyróżnienie treściowe & \wyroznienie{ważne} \\
\kod{\textbackslash osoba\{Nowak\}\{Jan\}} & osoba & Jan Nowak plus wpis do indeksu \\
\bottomrule
\end{tabular}
\end{tabelaepi}

Rozróżnienie kursywy i wytłuszczenia nie jest kwestią gustu: Instrukcja przewiduje kursywę dla wyrazów obcych i tytułów, a wytłuszczenie dla wyróżnień w tekście \parencite{instrukcjaISI}.

Polecenie \kod{\textbackslash osoba} stosuj \wyroznienie{od pierwszego dnia}. Indeks nazwisk powstaje z tych poleceń, a uzupełnianie go na końcu to kilka godzin pracy.

\section{Poziom drugi: tabele, rysunki, wzory}

\subsection{Listy i cytaty}

Listy stosuj tam, gdzie kolejność albo równorzędność elementów jest istotna. Nie zastępuj nimi wywodu -- praca złożona z wypunktowań czyta się jak prezentacja i tak jest oceniana.

Fragment dłuższy niż dwa lub trzy zdania wyodrębniasz graficznie, czego wymaga Instrukcja:

\begin{zapis}
\begin{quote}
Kryterium, w ktorym oceny wariantow sa do siebie zblizone, niesie malo
informacji rozniczujacej i powinno otrzymac wage nizsza.
\end{quote}
\end{zapis}

\subsection{Tabele}

Tabela w pracy ma tytuł nad sobą i źródło pod sobą. Robi to środowisko \kod{tabelaepi}, które przyjmuje tytuł i etykietę.

\begin{zapis}
\begin{tabelaepi}{Moduly pakietu}{tab:moduly}
\begin{tabular}{@{}lll@{}}
\toprule
Plik & Odpowiedzialnosc & Funkcje eksportowane \\
\midrule
\plik{algorytm.R} & obliczenia & \fun{oblicz_ranking} \\
\bottomrule
\end{tabular}
\zrodlo{oprac. wlasne}
\end{tabelaepi}
\end{zapis}

Wnętrze \kod{tabular} czyta się tak. Deklaracja kolumn podaje ich wyrównanie: \kod{l} do lewej, \kod{r} do prawej, \kod{c} do środka, \kod{p\{3cm\}} kolumna o stałej szerokości z łamaniem wiersza. Zapis \kod{@\{\}} na brzegach usuwa wcięcie. Komórki rozdziela ampersand, a wiersz kończą dwa ukośniki. Linie poziome dają \kod{\textbackslash toprule}, \kod{\textbackslash midrule} i \kod{\textbackslash bottomrule} -- i tylko one. Tabela naukowa nie ma linii pionowych ani kratownicy.

Tabela szersza niż strona: użyj \kod{tabularx} z kolumną \kod{X}, która sama dobierze szerokość. Tabela dłuższa niż strona przenosi się do aneksu, tak stanowi Instrukcja.

\subsection{Rysunki}

\begin{zapis}
\begin{rysunekepi}{Mapa decyzyjna wariantow}{rys:mapa}
\includegraphics[width=0.78\textwidth]{przyklad-wykres.png}
\zrodlo{oprac. wlasne}
\end{rysunekepi}
\end{zapis}

Pliki graficzne trzymaj w katalogu \plik{rysunki/} -- klasa szuka ich tam automatycznie, więc podajesz samą nazwę pliku. Szerokość podawaj względem szerokości tekstu, nie w centymetrach. Wykresy z R zapisuj w wysokiej rozdzielczości:

\begin{lstlisting}[numbers=none]
png("rysunki/wykres.png", width = 1600, height = 1000, res = 220)
plot(x, y)
dev.off()
\end{lstlisting}

Wykres liczbowy wstawiaj przez środowisko \kod{wykresepi}; ma osobną numerację, tak jak wymaga Instrukcja.

Tabele i rysunki są \termin{pływakami}: LaTeX umieszcza je tam, gdzie wychodzi najlepiej, niekoniecznie w miejscu wpisania. To nie jest usterka. Nie walcz z tym w trakcie pisania -- zajmij się tym przy redakcji, a w tekście zawsze odsyłaj przez etykietę, nigdy przez zwrot \enquote{poniższa tabela}.

\subsection{Wzory}

Wzory składa się w \termin{trybie matematycznym}, w którym symbole ustawiane są inaczej niż zwykły tekst: kursywą i z własnymi odstępami. Wzór w linii otacza się znakami dolara, które ten tryb włączają i wyłączają. Wzór wyróżniony i numerowany zapisuje się tak:

\begin{zapis}
\begin{equation}\label{eq:cc}
  CC_i = \frac{d_i^{-}}{d_i^{+} + d_i^{-}}.
\end{equation}

We wzorze \eqref{eq:cc} symbol $d_i^{+}$ oznacza...
\end{zapis}

Daje to wynik:

\begin{equation}\label{eq:cc}
  CC_i = \frac{d_i^{-}}{d_i^{+} + d_i^{-}}.
\end{equation}

Wzór bez numeru zapisuje się w nawiasach kwadratowych poprzedzonych ukośnikiem. Kilka wzorów wyrównanych do znaku równości składa środowisko \kod{align}.

\begin{tabelaepi}{Zapis matematyczny}{tab:matematyka}
\begin{tabular}{@{}ll@{}}
\toprule
Chcesz & Piszesz \\
\midrule
ułamek & \kod{\textbackslash frac\{a\}\{b\}} \\
indeks dolny i górny & \kod{x\_i}, \kod{x\textasciicircum 2}, \kod{x\_\{ij\}\textasciicircum\{+\}} \\
suma, iloczyn & \kod{\textbackslash sum\_\{i=1\}\textasciicircum\{n\}}, \kod{\textbackslash prod\_\{j=1\}\textasciicircum\{m\}} \\
pierwiastek & \kod{\textbackslash sqrt\{x\}} \\
nawiasy skalujące się & \kod{\textbackslash left( ... \textbackslash right)} \\
litery greckie & \kod{\textbackslash alpha}, \kod{\textbackslash lambda}, \kod{\textbackslash Phi} \\
wartość bezwzględna & \kod{\textbackslash lvert x \textbackslash rvert} \\
tekst wewnątrz wzoru & \kod{\textbackslash text\{dla \} x > 0} \\
macierz & \kod{\textbackslash begin\{pmatrix\} a \& b \textbackslash\textbackslash{} c \& d \textbackslash end\{pmatrix\}} \\
\bottomrule
\end{tabular}
\end{tabelaepi}

Symbole zostawiaj takie jak w artykule źródłowym -- czytelnik ma móc porównać. I pamiętaj o zasadzie z podrozdziału~\ref{sec:rozdzial1}: po każdym wzorze następuje akapit objaśniający, co ten wzór wyraża i po co jest w pracy.

\subsection{Listingi kodu}

\termin{Listing} to fragment kodu wydzielony z tekstu i złożony pismem maszynowym, z własnym numerem i podpisem, tak samo jak tabela czy rysunek. Standardy mówią o \termin{fragmentach kodu} \parencite{standardyEPI}; nazwa \enquote{listing} jest utrwaloną nazwą tego elementu w pracach informatycznych i tak też podpisuje go wzór, numerując listingi ciągle przez całą pracę.

\begin{zapis}
\begin{lstlisting}[caption={Wyznaczanie wag},label={lst:wagi}]
oblicz_wagi <- function(macierz) {
  stopifnot(is.matrix(macierz))
  colSums(macierz) / sum(macierz)
}
\end{lstlisting}
\end{zapis}

Kolorowanie składni, numerację wierszy i ramkę ustawia klasa. Kod z pliku wstawisz poleceniem \kod{\textbackslash listingR}, które przyjmuje podpis, etykietę i ścieżkę -- wtedy listing zawsze odpowiada aktualnej wersji kodu.

\wyroznienie{Wnętrze listingu musi być zapisane wyłącznie znakami ASCII.} Komentarze w kodzie pisz bez polskich znaków diakrytycznych. Powód jest podwójny. Tablica znaków pakietu składającego listingi obejmuje tylko ASCII, więc polskie litery nie tyle znikają, co \wyroznienie{wędrują na początek wyrazu}: zapisane w kodzie słowo \kod{ścieżka} złoży się jako \kod{śżcieka}. Usterka jest cicha, bo kompilacja kończy się bez błędu. Sprawdzanie pakietu R zgłasza z kolei znaki spoza ASCII w kodzie źródłowym jako problem przenośności, więc to samo ograniczenie obowiązuje w repozytorium. Podpis listingu jest zwykłym tekstem pracy i polskich znaków używać może; ograniczenie dotyczy wyłącznie wnętrza. Kontroli służy skrypt \plik{tools/sprawdz-listingi.R}, ale wymaga on R, więc działa tylko \wyroznienie{lokalnie}. \wyroznienie{Na Overleaf} objaw widać wprost w złożonym dokumencie: litery w listingu są poprzestawiane.

Zostaje pytanie, co zrobić, gdy trzeba pokazać tekst polski złożony pismem maszynowym: komunikat błędu, wynik działania programu, fragment danych. Do tego służy osobne środowisko \kod{wydruk}, które \wyroznienie{polskie znaki składa poprawnie}, bo korzysta z innego pakietu:

\begin{zapis}
\begin{wydruk}
Blad: kolumna "ocena jakosci" zawiera wartosci spoza zakresu.
\end{wydruk}
\end{zapis}

Cena jest taka, że nie ma tam kolorowania składni ani numeracji wierszy. Podział jest więc prosty: \wyroznienie{kod idzie do listingu, wyjście programu do wydruku}. Kodowi ograniczenie i tak nie przeszkadza, bo sprawdzanie pakietu R wymaga ASCII niezależnie od składu.

\section{Poziom trzeci: przy redakcji}

\subsection{Praca w wielu plikach}

Rozdziały są osobnymi plikami włączanymi poleceniem \kod{\textbackslash input}, podawanym bez rozszerzenia. Dzięki temu nie przewijasz tysiąca wierszy i łatwiej znaleźć miejsce błędu. Żeby przy pisaniu jednego rozdziału kompilować tylko jego, zakomentuj pozostałe polecenia włączenia w pliku głównym. Przed oddaniem odkomentuj wszystkie -- inaczej numeracja, spis treści i bibliografia będą niepełne.

Znak procentu wyłącza resztę wiersza. Komentarze służą do notatek dla siebie i znikają z gotowego pliku. Przed oddaniem przejrzyj je i usuń te, które są notatkami roboczymi.

\subsection{Kontrola typografii}

Klasa blokuje \termin{wdowy i sieroty}, czyli pojedyncze wiersze akapitu odcięte na końcu albo na początku strony, ale zdarza się, że LaTeX nie umie złamać zbyt długiego słowa i wypuszcza je poza margines. W dzienniku kompilacji zobaczysz wtedy komunikat \kod{Overfull \textbackslash hbox}. Naprawiaj w tej kolejności: przeformułuj zdanie, co najczęściej jest najlepszym rozwiązaniem; podpowiedz miejsce podziału zapisem \kod{wielo\textbackslash-kryterialny}; przy długim adresie sieciowym użyj polecenia \kod{\textbackslash url}, które pozwala łamać w sensownych miejscach. Liczba nadmiarowych pudełek w gotowej pracy powinna być zerowa albo bliska zeru.

\subsection{Sekwencja kompilacji}

Pełne złożenie pracy wymaga kilku przebiegów, bo spis treści, odsyłacze, bibliografia i indeks potrzebują danych z przebiegu poprzedniego. Kolejność to LuaLaTeX, potem \pkg{biber}, który składa bibliografię z pliku źródeł, potem \pkg{makeindex}, który porządkuje indeks nazwisk, a na koniec dwa razy LuaLaTeX.

\wyroznienie{Na Overleaf.} Sekwencja uruchamia się sama po naciśnięciu przycisku kompilacji i nie musisz o niej pamiętać.

\wyroznienie{Lokalnie.} Całą sekwencję uruchamia za Ciebie \pkg{latexmk}, wywoływany z katalogu pracy:

\begin{lstlisting}[language=bash,numbers=none]
latexmk -lualatex main.tex
\end{lstlisting}

Skutek praktyczny jest w obu ścieżkach ten sam: \wyroznienie{po dodaniu nowego cytowania albo etykiety pierwsza kompilacja może pokazać znaki zapytania} -- po drugiej znikną.

Przy okazji powstaje kilkanaście \termin{plików pomocniczych}, o rozszerzeniach \plik{.aux}, \plik{.toc} czy \plik{.bbl}. To w nich LaTeX przechowuje między przebiegami numery rozdziałów, strony odsyłaczy i złożoną bibliografię. Przydaje się to do momentu, w którym kompilacja zaczyna zachowywać się dziwnie, bo uszkodzony plik pomocniczy potrafi utrzymywać błąd, którego w źródle już nie ma. Czyszczenie wygląda wtedy różnie w obu ścieżkach: \wyroznienie{lokalnie} usuwa je polecenie \kod{latexmk -c}, \wyroznienie{na Overleaf} służy do tego \termin{Recompile from scratch} z rozwijanego menu obok przycisku kompilacji.

\section{Diagnostyka}\label{sec:diagnostyka}

\subsection{Jak czytać błąd}

Komunikat zaczyna się od wykrzyknika, a numer wiersza stoi w linii rozpoczynającej się od litery \kod{l} z kropką. Trzy rzeczy warto wiedzieć. Wiersz wskazany w błędzie bywa o jeden dalej niż przyczyna, bo brakujący nawias zamykający zgłasza się dopiero tam, gdzie coś przestało się zgadzać. Pierwszy błąd jest ważny, a reszta zwykle z niego wynika, więc napraw pierwszy i skompiluj ponownie, zamiast czytać wszystkie. Na Overleaf przełącz widok na pełny dziennik, gdy podsumowanie nie wystarcza.

\subsection{Katalog błędów}

\begin{tabelaepi}{Najczęstsze błędy kompilacji}{tab:bledy}
\begin{tabular}{@{}p{4.5cm}p{4.6cm}p{4.6cm}@{}}
\toprule
Komunikat albo objaw & Przyczyna & Naprawa \\
\midrule
\kod{Undefined control sequence} & literówka w poleceniu & sprawdź pisownię w tabeli~\ref{tab:makra} \\
\kod{Missing \$ inserted} & znak matematyczny poza trybem matematycznym & otocz znakami dolara albo poprzedź ukośnikiem \\
\kod{Missing \textbackslash begin\{document\}} & tekst wpisany w preambule & przenieś treść za początek dokumentu \\
\kod{File ... not found} & zła nazwa albo ścieżka grafiki & grafiki trzymaj w \plik{rysunki/} \\
\kod{Environment ... undefined} & literówka w nazwie środowiska & porównaj początek i koniec \\
\kod{\textbackslash begin ended by \textbackslash end} & niedomknięte środowiska & sprawdź parowanie \\
\kod{Extra alignment tab} & za dużo ampersandów w wierszu & policz kolumny w deklaracji \\
\kod{Misplaced alignment tab} & ampersand poza tabelą & poprzedź go ukośnikiem \\
\kod{Citation ... undefined} & klucz nieobecny w pliku bibliograficznym & sprawdź klucz, skompiluj dwa razy \\
\kod{Reference ... undefined} & odsyłacz do nieistniejącej etykiety & sprawdź etykietę, skompiluj dwa razy \\
\kod{Empty bibliography} & brak cytowań w tekście & dodaj cytowanie \\
\kod{Overfull \textbackslash hbox} & tekst wychodzi poza margines & przeformułuj zdanie \\
\kod{Underfull \textbackslash hbox} & zbyt rozstrzelony wiersz & zwykle można zignorować \\
\kod{There's no line here to end} & podwójny ukośnik bez wiersza & usuń go, do akapitu wystarczy pusta linia \\
\kod{Paragraph ended before ...} & brakujący nawias klamrowy & sprawdź parowanie \\
\kod{Package inputenc Error} & plik w innym kodowaniu niż UTF-8 & zapisz ponownie w UTF-8 \\
przestawione polskie litery w listingu & znaki spoza ASCII w kodzie & usuń diakrytykę \\
kod bez właściwego kroju pisma & dokument złożony pdfLaTeX-em & ustaw kompilator na LuaLaTeX \\
brak indeksu nazwisk & brak polecenia \kod{\textbackslash osoba} & dodaj polecenia i skompiluj \\
dokument nie odświeża się mimo zmian & uszkodzone pliki pomocnicze & pełne złożenie od zera; lokalnie \kod{latexmk -c} \\
\bottomrule
\end{tabular}
\end{tabelaepi}

\subsection{Gdy nic nie pomaga}

Wykonaj pełną kompilację od zera po usunięciu plików pomocniczych. Następnie zakomentuj połowę treści i skompiluj, żeby zawęzić, w której połowie jest problem; powtórz kilka razy, aż zostanie kilka wierszy. Na koniec sprawdź, czy problem występuje też w dokumencie pokazowym dołączonym do wzoru -- jeżeli tak, przyczyna leży w środowisku, a nie w Twoim tekście.

\chapter{Bibliografia: plik BibTeX i Zotero}\label{ch:bibliografia}

Bibliografia w pracy powstaje automatycznie z pliku \plik{bibliografia/literatura.bib}. Nie składasz jej ręcznie i nie sortujesz -- robi to styl odwzorowujący opisy z Instrukcji \parencite{instrukcjaISI}. Twoje zadanie sprowadza się do jednego: poprawnych wpisów w pliku. Opis wychodzi wtedy zgodny z wymogami sam.

Zasada, której pilnuje recenzent: bibliografia zawiera wyłącznie pozycje faktycznie wykorzystane, do których odsyłają powołania w tekście. Pozycje przeczytane, ale niewykorzystane, do bibliografii nie trafiają.

\wyroznienie{Prowadź plik bibliograficzny od pierwszego dnia.} Odtwarzanie bibliografii po napisaniu pracy zajmuje kilka dni i zawsze kończy się brakującym numerem strony.

\section{Anatomia wpisu}

\begin{zapis}
@book{cisek2002filozoficzne,
  author    = {Cisek, Sabina},
  year      = {2002},
  title     = {Filozoficzne aspekty informacji naukowej},
  location  = {Krakow},
  publisher = {Wydaw. UJ},
  langid    = {polish},
}
\end{zapis}

Pierwsza linia podaje typ wpisu, który decyduje o postaci opisu, oraz klucz cytowania, którym powołujesz się w tekście. Dalej idą pola; wartość zawsze w nawiasach klamrowych. Powołanie \kod{\textbackslash parencite\{cisek2002filozoficzne\}} daje w tekście \enquote{(Cisek 2002)}.

\subsection{Klucze}

Klucz jest Twój i nigdzie nie widać go w gotowym dokumencie, ale przy stu pozycjach zaczyna mieć znaczenie. Konwencja przyjęta w seminarium to nazwisko, rok i pierwsze znaczące słowo tytułu, bez znaków diakrytycznych i bez spacji: \kod{cisek2002filozoficzne}, \kod{callaway2024continuous}, \kod{liPearl2024probabilities}.

Klucz raz nadany zostaje na zawsze. Zmiana klucza po napisaniu połowy pracy oznacza przeszukanie wszystkich rozdziałów.

\subsection{Pola nazwisk}

Autorów rozdziela słowo \kod{and}, \wyroznienie{nigdy przecinek}. Zapis \kod{Nazwisko, Imie} jest jednoznaczny i jego się trzymaj. Przy nazwisku złożonym albo nazwie instytucji użyj podwójnych nawiasów, żeby całość została potraktowana jako jedna jednostka:

\begin{zapis}
author = {Nowak, Jan and Zielinski, Kazimierz}
author = {{Biblioteka Jagiellonska}}
author = {van Dijk, Teun A.}
\end{zapis}

Imiona podawaj w pełnym brzmieniu -- Instrukcja wymaga pełnych imion w opisie bibliograficznym. Inicjały w powołaniu styl doda sam, gdy będą potrzebne do odróżnienia dwóch autorów o tym samym nazwisku.

Polskie znaki wpisuj wprost, a plik zapisz w UTF-8. LaTeX zmienia wielkość liter w tytułach w części stylów, więc skrót albo nazwę własną otocz dodatkowymi nawiasami klamrowymi: \kod{title = \{Zastosowanie metody \{TOPSIS\}\}}.

\section{Typy wpisów}

Poniższe wpisy odpowiadają wprost przykładom opisu z Instrukcji. Komplet z komentarzami znajdziesz w pliku bibliograficznym dołączonym do wzoru -- zacznij od skopiowania właściwego wzorca i podmiany wartości.

\subsection{Książka i praca zbiorowa}

Książka wielu autorów różni się od pojedynczej tylko listą w polu \kod{author}. Praca zbiorowa ma redaktora zamiast autora, a styl sam dopisze skrót \enquote{red.}:

\begin{zapis}
@collection{zielinski1998swiat,
  editor    = {Zielinski, Jan},
  year      = {1998},
  title     = {Swiat komputerow},
  location  = {Wroclaw},
  publisher = {Globus},
  langid    = {polish},
}
\end{zapis}

Daje to opis \enquote{Zieliński, Jan red. (1998). \termin{Świat komputerów}. Wrocław: Globus.}, a w tekście powołanie \enquote{(Zieliński red. 1998)}.

\subsection{Rozdział w pracy zbiorowej}

\begin{zapis}
@incollection{nowak2004srodowisko,
  author    = {Nowak, Jan},
  year      = {2004},
  title     = {Srodowisko informacyjne},
  editor    = {Makowski, Eugeniusz},
  booktitle = {Czlowiek wspolczesny},
  location  = {Warszawa},
  publisher = {Wydaw. Atrakcja},
  pages     = {11-28},
  langid    = {polish},
}
\end{zapis}

\subsection{Artykuł}

W pozycjach polskich podajesz \wyroznienie{numer} w polu \kod{number}, w obcojęzycznych zwyczajowo \wyroznienie{tom} w polu \kod{volume}. Pole \kod{langid} przełącza skróty na oryginalne, dzięki czemu artykuł angielski dostaje oznaczenia \enquote{vol.} i \enquote{pp.} zamiast \enquote{nr} i \enquote{s.}:

\begin{zapis}
@article{hjorland1998theory,
  author       = {Hjorland, Birger},
  year         = {1998},
  title        = {Theory and metatheory of information science},
  journaltitle = {Journal of Documentation},
  volume       = {5},
  pages        = {606-621},
  langid       = {english},
}
\end{zapis}

W artykule prasowym datę dzienną wpisz w pole \kod{issue}; stanie wtedy między numerem a stronami, zgodnie z przykładem z Instrukcji.

\subsection{Hasło w wydawnictwie informacyjnym}

Hasło z własnym autorem opisuje się jak rozdział w pracy zbiorowej. Hasło bez autora wymaga podania autora dzieła nadrzędnego w polu \kod{bookauthor}:

\begin{zapis}
@incollection{apoteoza1996,
  title      = {Apoteoza},
  bookauthor = {Kopalinski, Wladyslaw},
  year       = {1996},
  booktitle  = {Slownik wyrazow obcych},
  location   = {Warszawa},
  publisher  = {Wiedza Powszechna},
  pages      = {14},
  langid     = {polish},
}
\end{zapis}

\subsection{Dokumenty sieciowe}

Do wpisu drukowanego dodajesz pola \kod{doi}, \kod{url} i \kod{urldate}. Datę odczytu zapisuj w formacie rok-miesiąc-dzień; styl przestawi ją na polską postać z kropkami. Strony internetowe i wpisy w blogach opisuje typ \kod{@online}, a pole \kod{keywords} o wartości \kod{netografia} pozwala złożyć te pozycje w odrębnym wykazie, jeżeli uzgodnisz to z promotorem.

\subsection{Pozycje szczególne}

Dokument na nośniku materialnym opisuje typ nośnika w polu \kod{howpublished}. Norma wymaga pól \kod{entrysubtype} o wartości \kod{norma} oraz \kod{shorthand}, bo rok wchodzi w skład jej oznaczenia i nie powtarza się go w nawiasie \parencite{pniso690}. Pozycja bez autora potrzebuje pola \kod{shorttitle} z dwoma pierwszymi słowami tytułu, żeby powołanie miało postać wymaganą przez Instrukcję.

\section{Zestawienie pól}

\begin{tabelaepi}{Pola wpisu bibliograficznego}{tab:pola}
\begin{tabular}{@{}llp{6.6cm}@{}}
\toprule
Pole & Kiedy & Uwaga \\
\midrule
\kod{author} lub \kod{editor} & zawsze, gdy są & zapis nazwisko, imię; rozdzielone słowem \kod{and} \\
\kod{year} & zawsze & sam rok; brak roku wolno pominąć \\
\kod{title} & zawsze & \\
\kod{shorttitle} & pozycje bez autora & dwa pierwsze słowa tytułu \\
\kod{journaltitle} & artykuł & pełny tytuł, bez skrótów \\
\kod{booktitle} & rozdział, hasło & tytuł dzieła nadrzędnego \\
\kod{bookauthor} & hasło bez autora & autor dzieła nadrzędnego \\
\kod{number} & pozycje polskie & daje oznaczenie \enquote{nr 2} \\
\kod{volume} & pozycje obcojęzyczne & daje oznaczenie \enquote{vol. 5} \\
\kod{issue} & gazeta & data dzienna \\
\kod{pages} & artykuł, rozdział & zakres zapisany łącznikiem \\
\kod{location} & wydawnictwo zwarte & miasto \\
\kod{publisher} & wydawnictwo zwarte & \\
\kod{doi} & gdy jest & sam identyfikator, bez adresu \\
\kod{url} & dokument sieciowy & \\
\kod{urldate} & zawsze przy adresie & format rok-miesiąc-dzień \\
\kod{langid} & zawsze & \kod{polish} albo \kod{english} \\
\kod{keywords} & opcjonalnie & wartość \kod{netografia} \\
\kod{howpublished} & nośnik materialny & na przykład płyta DVD \\
\kod{entrysubtype} & norma & wartość \kod{norma} \\
\kod{shorthand} & norma & oznaczenie do powołania \\
\bottomrule
\end{tabular}
\end{tabelaepi}

Pole \kod{langid} decyduje o skrótach w opisie, a jego brak jest traktowany jak pozycja polska. \wyroznienie{Wypełniaj je zawsze} -- inaczej artykuł angielski dostanie polskie oznaczenia tomu i stron.

\section{Zotero}

Ręczne pisanie stu wpisów jest wykonalne, ale nierozsądne. Zotero zbiera metadane z przeglądarki i eksportuje je do pliku bibliograficznego.

\subsection{Konfiguracja}

Zainstaluj Zotero wraz z wtyczką do przeglądarki, a następnie dodatek Better BibTeX. Dodatek jest niezbędny: daje stabilne klucze cytowania, eksport w formacie biblatex i automatyczne odświeżanie pliku. Utwórz kolekcję dla pracy, a w niej podkolekcje na metodę, dziedzinę i narzędzia.

W ustawieniach dodatku ustaw wzorzec klucza odpowiadający konwencji seminaryjnej i \wyroznienie{włącz pinowanie kluczy}. Bez tego klucz może się zmienić po edycji metadanych i powołania w pracy przestaną działać.

Eksport ustawia się raz: menu podręczne kolekcji, polecenie eksportu, format \termin{Better BibLaTeX}, opcja utrzymywania pliku w aktualności. Jako cel wskaż plik bibliograficzny w projekcie pracy. Wybieraj format biblatex, a nie starszy BibTeX -- styl korzysta z pól, których starszy format nie zna.

\subsection{Odpowiedniki typów}

\begin{tabelaepi}{Typy dokumentów w Zotero i we wpisie bibliograficznym}{tab:typyzotero}
\begin{tabular}{@{}ll@{}}
\toprule
Typ w Zotero & Typ wpisu \\
\midrule
Book & \kod{@book} \\
Book Section & \kod{@incollection} \\
Journal Article & \kod{@article} \\
Magazine, Newspaper Article & \kod{@article} \\
Conference Paper & \kod{@inproceedings} \\
Thesis & \kod{@thesis} \\
Report & \kod{@report} \\
Web Page, Blog Post & \kod{@online} \\
Computer Program & \kod{@software} \\
Dictionary, Encyclopedia Entry & \kod{@incollection} \\
\bottomrule
\end{tabular}
\end{tabelaepi}

\subsection{Poprawki po imporcie}

Metadane z baz bibliograficznych bywają niekompletne. Po każdym imporcie sprawdź imiona, bo bazy notorycznie zostawiają inicjały, a Instrukcja wymaga pełnych imion. Uzupełnij pole języka wartością \kod{polish} albo \kod{english}, bo stąd bierze się \kod{langid}, i po eksporcie zajrzyj do pliku, czy pole rzeczywiście przeszło. Rozwiń skrócone tytuły czasopism. Sprawdź typ dokumentu, bo rozdział zaimportowany jako artykuł da zły opis. Uzupełnij miejsce wydania, które Zotero często gubi przy książkach. Usuń przedrostek adresowy z identyfikatora cyfrowego. Zapisz tytuł tak, jak w oryginale publikacji, bo bazy anglojęzyczne stosują kapitaliki wyrazowe.

Zasada praktyczna: \wyroznienie{nigdy nie ufaj pierwszemu importowi}. Sprawdzenie wpisu zajmuje minutę, a poprawianie bibliografii tydzień przed obroną -- znacznie więcej.

\section{Powołania w tekście}

\begin{tabelaepi}{Postacie powołania}{tab:powolania}
\begin{tabular}{@{}llp{5cm}@{}}
\toprule
Sytuacja & Polecenie & Efekt \\
\midrule
zwykłe & \kod{\textbackslash parencite\{klucz\}} & (Cisek 2002) \\
ze stroną & \kod{\textbackslash parencite[s.\textasciitilde 15-21]\{klucz\}} & (Kowalski 2010a, s. 15--21) \\
kilka pozycji & \kod{\textbackslash parencite\{k1,k2\}} & (Cisek 2002; Nowak 2006) \\
w zdaniu & \kod{\textbackslash textcite\{klucz\}} & Woźniak (1997) pokazuje\ldots \\
z przedrostkiem & \kod{\textbackslash parencite[zob.][s.\textasciitilde 9]\{klucz\}} & (zob. Nowak 2006, s. 9) \\
\bottomrule
\end{tabular}
\end{tabelaepi}

Inicjał przy zbieżnych nazwiskach, skrót przy więcej niż trzech autorach i sufiksy rocznika styl dobiera automatycznie. Nie wpisuj ich ręcznie w pliku bibliograficznym.

Cytat dosłowny z tłumaczeniem własnym zapisujesz razem z oznaczeniem tłumacza i oryginałem w przypisie dolnym:

\begin{zapis}
\enquote{Przetlumaczony fragment} [tlum. wlasne --- JK]
\parencite[s.~44]{hjorland1998theory}\footnote{\enquote{Original wording}.}
\end{zapis}

\section{Cytowanie oprogramowania i danych}

Pakiety R i zbiory danych trafiają do \wyroznienie{wykazu źródeł}, a nie do bibliografii; rozróżnienie omawia podrozdział~\ref{sec:wykazzrodel}. Jeżeli jednak powołujesz się na pakiet w tekście jak na publikację, potrzebujesz wpisu. Dane pobierzesz w konsoli:

\begin{lstlisting}[numbers=none]
citation("ggplot2")
toBibtex(citation("ggplot2"))
\end{lstlisting}

Twój pakiet powinien mieć plik \plik{inst/CITATION}, żeby polecenie \fun{citation} zwracało poprawny opis. Po nadaniu wydania i uzyskaniu identyfikatora cyfrowego opis własnego pakietu wygląda tak:

\begin{zapis}
@software{kowalski2027nazwapakietu,
  author  = {Kowalski, Jan},
  year    = {2027},
  title   = {NazwaPakietu: krotki opis przeznaczenia},
  version = {1.0.0},
  doi     = {10.5281/zenodo.0000000},
  url     = {https://github.com/nazwauzytkownika/NazwaPakietu},
  urldate = {2027-05-15},
  langid  = {polish},
}
\end{zapis}

Zbiór danych opisujesz jak dokument sieciowy, z podaniem wersji albo daty pobrania oraz licencji w polu \kod{note}.

\section{Kontrola}

Przed oddaniem sprawdź trzy rzeczy. Po pierwsze zgodność w obie strony: każda pozycja bibliografii ma powołanie w tekście, a każde powołanie pozycję w bibliografii. Nierozwiązane powołania widać w dzienniku kompilacji; pozycji bez powołania styl po prostu nie złoży, co działa na Twoją korzyść. Po drugie obecność pozycji obcojęzycznych, wymaganych zarówno przez Standardy, jak i przez Instrukcję. Po trzecie kompletność opisów -- przejrzyj gotową bibliografię w złożonym dokumencie, a nie w pliku źródłowym, bo brakujące miasto, ucięty tytuł czasopisma albo brak daty odczytu widać dopiero w składzie.

Gdy bibliografia jest pusta albo nieaktualna po zmianach w pliku, uruchom pełną kompilację od zera.

\section{Wariant APA}

Jeżeli uzgodnisz z promotorem styl APA zamiast instytutowego, wystarczy dopisać opcję klasy \kod{apa} w poleceniu \kod{\textbackslash documentclass}. Ten sam plik bibliograficzny obsługuje oba style. Łączniki są polskie, a nazwy miesięcy w datach dostępu pozostają w formie właściwej dla angielskiego zapisu APA.

Wybór stylu uzgadnia się \wyroznienie{raz, na początku} -- Instrukcja wprost przewiduje ustalenie preferowanego stylu z promotorem. Zmiana w trakcie pisania oznacza przegląd wszystkich wpisów pod kątem pól wymaganych przez nowy styl.

\chapter{Overleaf i współpraca}\label{ch:overleaf}

Overleaf to LaTeX w przeglądarce. Nie trzeba nic instalować, projekt jest dostępny z każdego komputera, a promotor widzi tę samą wersję co Ty. Dla pracy licencjackiej pisanej pod opieką promotora to najmniej kłopotliwe środowisko. Praca lokalna też jest możliwa i opisuje ją podrozdział~\ref{sec:lokalnie}.

\section{Ścieżka studenta}\label{sec:overleafstudent}

\subsection{Konto i projekt}

Załóż konto \wyroznienie{na adres uczelniany}. Sprawdź, czy Uniwersytet udostępnia licencję instytucjonalną -- jeżeli tak, konto w domenie uczelni otrzymuje funkcje płatne. Cała ścieżka opisana niżej działa również na planie bezpłatnym.

Wzór pobierasz z repozytorium seminarium: pobierz archiwum całego repozytorium, rozpakuj je, spakuj \wyroznienie{sam katalog wzoru} do nowego archiwum i wgraj je jako nowy projekt. Alternatywnie promotor udostępni gotowy projekt wzorcowy łączem tylko do odczytu; wtedy wystarczy skopiować projekt do własnego konta.

\subsection{Ustawienia}

\begin{tabelaepi}{Ustawienia projektu na Overleaf}{tab:overleaf}
\begin{tabular}{@{}llp{6.4cm}@{}}
\toprule
Ustawienie & Wartość & Dlaczego \\
\midrule
Compiler & LuaLaTeX & kod składa się wtedy właściwym krojem pisma \\
Main document & \plik{main.tex} & inaczej kompiluje się dokument pokazowy \\
Spell check & polski & podkreśla literówki w trakcie pisania \\
\bottomrule
\end{tabular}
\end{tabelaepi}

Ustawienie kompilatora jest jednorazowe i \wyroznienie{łatwe do przeoczenia}. Objaw pominięcia: praca kompiluje się poprawnie, ale kod i nazwy funkcji mają inny krój pisma niż powinny.

\subsection{Co jest czym w projekcie}

W pliku głównym uzupełniasz metadane na górze i nic więcej. Piszesz w plikach rozdziałów. Wpisy bibliograficzne trafiają do pliku w katalogu bibliografii, a grafiki do katalogu rysunków. Dokument pokazowy służy do zaglądania, gdy nie pamiętasz zapisu.

Klasy dokumentu ani plików stylu bibliograficznego \wyroznienie{nie ruszasz}. Zmiana czegokolwiek w nich oznacza, że praca przestaje odpowiadać wymogom Instytutu, a odpowiedzialność za to spada na Ciebie, nie na wzór. Jeżeli czegoś nie da się zrobić dostępnymi poleceniami, zgłoś to promotorowi -- brakuje wtedy polecenia we wzorze.

\subsection{Kompilacja i kopie}

Pierwsza kompilacja po dodaniu nowego powołania albo etykiety może pokazać znaki zapytania; po drugiej znikną. Rozwijane menu obok przycisku kompilacji zawiera polecenie pełnego przebiegu od zera -- użyj go, gdy dokument przestaje odpowiadać zmianom albo gdy bibliografia nie chce się odświeżyć.

Plan bezpłatny ma limit czasu kompilacji. Praca licencjacka mieści się w nim spokojnie, ale gdy zaczniesz się o niego ocierać, zakomentuj w pliku głównym włączenia rozdziałów, nad którymi akurat nie pracujesz. Przed oddaniem odkomentuj wszystkie.

Overleaf przechowuje projekt na swoich serwerach, ale \wyroznienie{to nie jest kopia zapasowa}. Raz w tygodniu pobierz źródła projektu i zapisz je poza komputerem, na którym pracujesz; przed każdym większym przemeblowaniem pracy dodatkowo. Historia zmian na planie bezpłatnym jest ograniczona i kopii zapasowej nie zastępuje.

\subsection{Udostępnienie i oddanie}

Udostępnij projekt promotorowi na jego adres, z prawem edycji. Na planie bezpłatnym każdy projekt można udostępnić jednej osobie -- dokładnie tyle, ile trzeba. Recenzentowi, gdy zajdzie potrzeba, udostępnij łącze tylko do odczytu; łącza nie liczą się do limitu współpracowników.

Do Archiwum Prac trafia plik wynikowy pobrany z projektu. Przed pobraniem wykonaj pełną kompilację od zera, sprawdź brak nierozwiązanych powołań i odsyłaczy oraz przejdź listę kontrolną z podrozdziału~\ref{sec:listakontrolna}.

\section{Plik bibliograficzny a Zotero}

Dodatek Better BibTeX odświeża plik bibliograficzny \wyroznienie{na Twoim dysku}. Overleaf o tym nie wie, więc plik trzeba do projektu przenieść. Najprostszy sposób to wgranie ręczne: w drzewie plików wskaż katalog bibliografii, prześlij zaktualizowany plik i potwierdź nadpisanie. Wystarczająco dobre, jeżeli robisz to raz na tydzień. Overleaf potrafi też wczytać publiczną bibliotekę Zotero i odświeżać ją na żądanie, co jest wygodne, ale wymaga ustawienia biblioteki jako publicznej. Przy dostępie do synchronizacji z repozytorium plik odświeża się razem z resztą projektu.

Niezależnie od sposobu: \wyroznienie{po każdej podmianie pliku uruchom pełną kompilację od zera}. Bibliografia budowana jest z pliku pomocniczego, który potrafi zostać na starej wersji.

\section{Ścieżka promotora}

\subsection{Rozdanie wzoru grupie}

Najmniej kłopotliwy układ, działający również na kontach bezpłatnych, wygląda tak. Prowadzący zakłada \wyroznienie{jeden projekt wzorcowy} z zawartości katalogu wzoru i włącza w nim łącze tylko do odczytu. Rozdaje to jedno łącze wszystkim uczestnikom. Każda osoba kopiuje projekt do własnego konta, otrzymując niezależną kopię, i udostępnia ją prowadzącemu z prawem edycji.

Zalety układu: łącze do odczytu nie zużywa limitu współpracowników, poprawka we wzorcu nie psuje niczyjej pracy, a prowadzący ma dostęp do wszystkich projektów z jednego pulpitu. Wadą jest to, że poprawka we wzorcu nie trafia do kopii już wykonanych -- dlatego \wyroznienie{wzór zamyka się przed pierwszymi zajęciami}, a późniejsze zmiany rozsyła jako opis, co podmienić.

\subsection{Przegląd prac}

Bez funkcji płatnych do dyspozycji pozostaje zwykła edycja i komentarze w treści. Przy kilkunastu pracach warto ustalić jedną konwencję i trzymać się jej przez cały rok. Konwencja proponowana w seminarium to uwagi zapisane jako komentarze LaTeX-a, które \wyroznienie{nie pojawiają się w gotowym dokumencie}:

\begin{zapis}
%% UWAGA: teza z tego akapitu nie wynika z przywolanego zrodla.
%% UWAGA: brakuje powolania.
%% PYTANIE: skad wartosc progu 0,7?
\end{zapis}

Uwaga stoi wtedy dokładnie przy problematycznym zdaniu, student widzi ją przy pisaniu, a gotowy dokument pozostaje czysty. Wyszukanie wszystkich uwag w projekcie to jedno kliknięcie w polu wyszukiwania.

Zasada zamykająca: \wyroznienie{student usuwa komentarz dopiero po naniesieniu poprawki}. Pusty wynik wyszukiwania oznacza, że wszystkie uwagi zostały zaadresowane; ten punkt znajduje się na liście kontrolnej przed oddaniem.

Przy licencji instytucjonalnej dochodzą śledzenie zmian, komentarze w panelu bocznym i pełna historia wersji -- wtedy powyższa konwencja staje się zbędna.

\subsection{Kontrola formalna}

Kontrolę formalną prowadź \wyroznienie{na złożonym dokumencie, nie w źródle}. Marginesy, interlinia, numeracja stron i układ elementów końcowych widać dopiero w składzie. Zestawienie wymogów wraz ze wskazaniem, co realizuje klasa, a co pozostaje po stronie autora, zawiera rozdział~\ref{ch:wymogi}.

W źródle warto sprawdzić dwie rzeczy: czy student nie zmieniał klasy dokumentu ani plików stylu bibliograficznego, co załatwia porównanie z wzorcem, oraz czy w projekcie nie zostały komentarze robocze i uwagi.

Przy kilkunastu pracach przegląd na żądanie nie działa. Sprawdzony układ to stałe okna: oddanie fragmentu do ustalonej daty, przegląd w ciągu tygodnia, omówienie na zajęciach. Fragmenty oddawane w kolejności pisania, a nie w kolejności rozdziałów -- kolejność podaje tabela~\ref{tab:kolejnosc}.

\section{Praca lokalna w Positronie}\label{sec:lokalnie}

Skoro pakiet R budujesz w Positronie, możesz w nim też pisać pracę. Kod i tekst są wtedy w jednym środowisku, a wykresy trafiają prosto do katalogu rysunków. Potrzebujesz dystrybucji TeX-a oraz rozszerzenia do LaTeX-a.

\begin{lstlisting}[language=bash,numbers=none]
cd wzor-pracy
latexmk -lualatex main.tex
\end{lstlisting}

Zaletami są brak limitu czasu kompilacji, praca bez dostępu do sieci i pełna historia w repozytorium. Wadami -- konieczność zainstalowania i utrzymania dystrybucji TeX-a oraz to, że udostępnienie promotorowi wymaga repozytorium zamiast jednego kliknięcia. Układ mieszany bywa najwygodniejszy: pakiet i wykresy lokalnie, tekst na Overleaf.

Repozytorium pakietu i pracę trzymaj \wyroznienie{osobno}. Praca zawiera pliki, które nie mają nic wspólnego z pakietem, a repozytorium pakietu ma być czyste -- to ono jest oceniane jako projekt dyplomowy i to jego adres podajesz na stronie tytułowej. Jeżeli prowadzisz pracę w repozytorium, ignoruj artefakty kompilacji.

\section{Problemy specyficzne dla Overleaf}

\begin{tabelaepi}{Problemy w pracy na Overleaf}{tab:problemyoverleaf}
\begin{tabular}{@{}p{5cm}p{4.5cm}p{4.5cm}@{}}
\toprule
Objaw & Przyczyna & Naprawa \\
\midrule
kod bez właściwego kroju pisma & kompilator ustawiony na pdfLaTeX & zmień kompilator na LuaLaTeX \\
kompiluje się dokument pokazowy & zły dokument główny & wskaż plik główny pracy \\
bibliografia nie odświeża się & stary plik pomocniczy & pełna kompilacja od zera \\
kompilacja przerwana limitem czasu & za dużo naraz przy szkicowaniu & zakomentuj nieużywane włączenia \\
grafika nie znaleziona & plik poza katalogiem rysunków & przenieś plik, podawaj samą nazwę \\
polskie znaki jako znaki zastępcze & plik w innym kodowaniu & zapisz ponownie w UTF-8 \\
zmiany współpracownika niewidoczne & stara karta przeglądarki & odśwież stronę \\
projekt nie chce się skopiować & łącze bez prawa kopiowania & poproś o łącze do odczytu \\
\bottomrule
\end{tabular}
\end{tabelaepi}

Katalog błędów kompilacji niezależnych od Overleaf zawiera podrozdział~\ref{sec:diagnostyka}.

\chapter{Jak napisać pracę}\label{ch:jakpisac}

Rozdział prowadzi przez wzór część po części. Każdy podrozdział odpowiada jednemu plikowi w katalogu rozdziałów. Kolejność podrozdziałów jest kolejnością w pracy, ale \wyroznienie{nie kolejnością pisania} --- o tym w podrozdziale~\ref{sec:kolejnosc}.

Twoja praca opisuje projekt: pakiet języka R implementujący algorytm z artykułu metodycznego. To zmienia sens niektórych części w stosunku do prac czysto teoretycznych. Stanem badań jest tu nie tylko literatura przedmiotu, ale też istniejące oprogramowanie. Materiałem badawczym są dane, na których pokazujesz działanie narzędzia. Wynikiem jest działające, sprawdzone narzędzie oraz wnioski z jego zastosowania.

\section{Wprowadzenie}

Plik \plik{rozdzialy/00-wprowadzenie.tex}, objętość trzy do pięciu stron, bez numeru rozdziału.

Wprowadzenie odpowiada czytelnikowi na cztery pytania w tej kolejności: o czym to jest, dlaczego to problem, co konkretnie zrobiłeś, jak czytać dalej. Recenzent czyta wprowadzenie i podsumowanie uważniej niż resztę --- z nich wyrabia sobie zdanie o pracy.

\subsection{Co musi zawierać}

\wyroznienie{Tło problemu.} Zacznij od problemu, który istnieje niezależnie od Twojego pakietu. Nie od narzędzia, tylko od pytania, na które badacz nauk społecznych dziś nie potrafi odpowiedzieć albo odpowiada z trudem. Dopiero na tym tle pojawia się metoda.

\wyroznienie{Luka.} Wyjaśnij, czego brakuje. W pracy projektowej luka ma dwie warstwy i warto rozróżnić je wprost. \termin{Luka technologiczna} to sytuacja, w której metoda jest opisana w literaturze, ale nie ma jej implementacji; implementacja istnieje, lecz wymaga przygotowania programistycznego; istniejące narzędzie nie obsługuje danych określonego rodzaju. \termin{Luka poznawcza} to z kolei to, czego badacze nie wiedzą, bo nie mają czym policzyć. Praca, która wypełnia wyłącznie pierwszą, jest sprawozdaniem z programowania, a nie pracą naukową.

Luka musi być sprawdzalna. Wymień narzędzia, które przejrzałeś, i napisz, czego w nich nie ma. Zdanie o nieistnieniu implementacji, podane bez wykazu przeszukanych źródeł, jest zarzutem pod Twoim adresem, a nie argumentem.

\wyroznienie{Cel.} Jedno zdanie w formie czynności: celem pracy jest zaprojektowanie, implementacja i weryfikacja pakietu języka R udostępniającego określoną procedurę badaczom nauk społecznych. Cel jest tym, co zrobiłeś, a nie tym, czego się dowiedziałeś.

\wyroznienie{Pytanie badawcze.} Cel mówi, co zbudowałeś; pytanie mówi, czego się dzięki temu dowiadujemy. Pytania zaczynające się od \enquote{jak} i \enquote{dlaczego} prowadzą do analizy, a \enquote{co} i \enquote{ile} do opisu. Sposób dochodzenia do pytania opisuje podrozdział~\ref{sec:pytanie}. Przykłady pytań właściwych dla tego typu pracy:

\begin{itemize}
\item Jakie warunki musi spełnić implementacja metody, żeby jej wyniki dały się odtworzyć przez osobę trzecią?
\item Które założenia metody okazują się najbardziej wrażliwe na jakość danych spotykanych w badaniach społecznych?
\item Jak wybór procedury generowania danych testowych wpływa na to, co wolno wywnioskować o poprawności implementacji?
\end{itemize}

\wyroznienie{Materiał i metoda.} Dwa lub trzy zdania: na jakich danych sprawdzasz narzędzie i jak weryfikujesz poprawność. Szczegóły idą do rozdziałów drugiego i trzeciego pracy.

\wyroznienie{Zakres i ograniczenia.} Co praca obejmuje, a czego nie. Ograniczenie postawione samodzielnie jest mocniejsze niż ten sam zarzut w recenzji.

\wyroznienie{Struktura.} Jedno zdanie o każdym rozdziale.

\subsection{Typowe błędy}

\begin{tabelaepi}{Błędy we wprowadzeniu}{tab:bledywprowadzenie}
\begin{tabular}{@{}p{6.6cm}p{6.6cm}@{}}
\toprule
Błąd & Dlaczego szkodzi \\
\midrule
Wprowadzenie zaczyna się od opisu pakietu & czytelnik nie wie jeszcze, po co mu ten pakiet \\
Cel sformułowany jako przedstawienie zagadnienia & nie da się sprawdzić, czy został osiągnięty \\
Luka stwierdzona bez wykazu przejrzanych narzędzi & twierdzenie nieudokumentowane \\
Historia dziedziny od połowy zeszłego wieku & wprowadzenie ma wprowadzać do problemu, nie do dyscypliny \\
Pytanie badawcze utożsamione z celem & brakuje warstwy poznawczej pracy \\
\bottomrule
\end{tabular}
\end{tabelaepi}

\begin{kontrolna}
\item Problem postawiony przed narzędziem
\item Luka wskazana i udokumentowana wykazem przejrzanych rozwiązań
\item Cel w jednym zdaniu, w formie czynności
\item Pytanie badawcze różne od celu
\item Zakres i ograniczenia wypisane
\item Zapowiedź struktury zgodna z faktyczną strukturą pracy
\item Każde twierdzenie o stanie rzeczy ma powołanie
\end{kontrolna}

\section{Rozdział pierwszy: podstawy metodyczne}\label{sec:rozdzial1}

Plik \plik{rozdzialy/01-podstawy.tex}, objętość osiem do dwunastu stron.

Odpowiednik sekcji poświęconej podstawom metodycznym w czasopismach o oprogramowaniu badawczym. Opisujesz matematykę i logikę, które zaimplementowałeś --- jeszcze nie kod. Rozdział ma udowodnić, że rozumiesz metodę, którą zamknąłeś w narzędziu.

\subsection{Zasada nadrzędna}

\wyroznienie{Po każdym wzorze następuje akapit wyjaśniający, co ten wzór robi i po co, językiem zrozumiałym dla badacza bez przygotowania matematycznego.} Wzór bez wyjaśnienia jest w tej pracy błędem merytorycznym, a nie oszczędnością miejsca. Twoim odbiorcą jest osoba, która chce metody użyć, a nie ją wyprowadzić.

Sprawdzian: przeczytaj sam akapit, bez wzoru. Czy nadal wiadomo, co się dzieje? Jeżeli nie, akapit jest za słaby.

\subsection{Co musi zawierać}

\wyroznienie{Osadzenie metody.} Jaki problem metoda rozwiązuje, skąd się wzięła, czym różni się od podejść wcześniejszych. Powołanie na artykuł źródłowy i dwa lub trzy opracowania wprowadzające.

\wyroznienie{Aparat formalny.} Wielkości wejściowe wraz z dziedzinami, kolejne kroki algorytmu, wielkości wyjściowe. Numeruj i etykietuj te wzory, do których praca się odwołuje. Objaśnij \wyroznienie{każdy} symbol, którego używasz --- nawet ten, który wydaje Ci się oczywisty.

\wyroznienie{Założenia i granice stosowalności.} Przy jakich warunkach metoda jest poprawna: rozkłady, niezależność obserwacji, kompletność danych, skale pomiarowe. I co się dzieje, gdy założenie jest naruszone. Ta sekcja zasila później kontrakt w specyfikacji oraz testy odporności pakietu --- pisząc ją, projektujesz jednocześnie oprogramowanie.

\wyroznienie{Rozwiązania pokrewne.} Przegląd istniejących implementacji w układzie problemowym: co już jest, czego brakuje, jak Twoje narzędzie się do tego ma. Nie lista, tylko wywód. Zakończ jednoznacznym stwierdzeniem, jaki wkład wnosi Twój pakiet.

\subsection{Skąd brać treść}

Z artykułu źródłowego, z sekcji definicyjnej i algorytmicznej. \wyroznienie{Nie przepisuj wzorów z opracowań wtórnych ani z zestawień tematycznych} --- zapis bywa tam skrócony i niedokładny. Wracaj do publikacji.

Wzoru nie tłumacz na polski w sensie notacji. Symbole zostają takie jak w źródle, żeby czytelnik mógł porównać; tłumaczysz sens, a nie znaki.

\begin{tabelaepi}{Błędy w rozdziale metodycznym}{tab:bledymetodyka}
\begin{tabular}{@{}p{6.6cm}p{6.6cm}@{}}
\toprule
Błąd & Dlaczego szkodzi \\
\midrule
Wzory przepisane bez objaśnienia symboli & recenzent nie odróżni zrozumienia od przepisania \\
Opis kodu zamiast opisu metody & kod należy do rozdziału drugiego \\
Pominięcie założeń metody & testy odporności nie mają z czego wyniknąć \\
Przegląd narzędzi w formie listy & brak wywodu problemowego, wymóg Standardów \\
Wzory przepisane z opracowania wtórnego & ryzyko powielenia cudzego błędu \\
\bottomrule
\end{tabular}
\end{tabelaepi}

\begin{kontrolna}
\item Każdy symbol objaśniony w tekście
\item Po każdym wzorze akapit o tym, co wzór robi
\item Założenia metody wypisane wprost
\item Granice stosowalności opisane
\item Przegląd rozwiązań pokrewnych w układzie problemowym
\item Wkład własny sformułowany jednoznacznie
\item Wzory sprawdzone z publikacją źródłową, nie z zestawieniem
\end{kontrolna}

\section{Rozdział drugi: implementacja i architektura}

Plik \plik{rozdzialy/02-implementacja.tex}, objętość dziesięć do czternastu stron.

Odpowiednik sekcji opisującej oprogramowanie. Najważniejszy rozdział inżynierski --- tu pokazujesz własny wkład. Standardy wymagają w tym miejscu opisu logiki aplikacji ze schematami oraz opisu implementacji z formatem danych i sposobem realizacji każdego modułu \parencite{standardyEPI}.

\subsection{Co musi zawierać}

\wyroznienie{Struktura projektu.} Katalogi, moduły, przepływ danych. Wyjaśnij, co gdzie leży i dlaczego akurat tak. Tabela modułów z odpowiedzialnościami i funkcjami eksportowanymi mówi więcej niż akapit opisu.

\wyroznienie{Przetwarzanie i walidacja danych.} Droga od danych surowych do struktury, na której pracuje algorytm: sprawdzanie warunków wstępnych, obsługa braków i kodów błędnych, skalowanie, agregacja. Wypisz, które naruszenia zatrzymują wykonanie i z jakim komunikatem --- to jest kontrakt Twojego narzędzia i najmocniejsza część tego rozdziału.

\wyroznienie{Silnik obliczeniowy.} Jak wzory z rozdziału pierwszego przełożyły się na kod. Pokaż fragment, nie całość; całość jest w repozytorium. Wyjaśnij decyzje, które nie wynikają wprost ze wzoru: sposób radzenia sobie z dzieleniem przez zero, wybór biblioteki optymalizacyjnej, zamianę pętli na operacje na wektorach.

\wyroznienie{System klas i metody generyczne.} Dlaczego wynik jest obiektem określonej klasy, a nie zwykłą ramką danych, i co dzięki temu potrafią funkcje wypisujące, podsumowujące i rysujące.

\wyroznienie{Procedura generowania danych testowych.} Rozkłady, parametry, zależności między zmiennymi, kontrolowane braki i wartości odstające, ziarno losowości. To jest element metodologii, a nie dodatek techniczny: dzięki znanej strukturze danych wiadomo, jaki wynik narzędzie \wyroznienie{powinno} zwrócić, więc odchylenie jest wykrywalne.

\wyroznienie{Weryfikacja poprawności.} Trzy warstwy, opisane osobno. Przypadki analityczne o wyniku policzonym ręcznie: mała macierz, przypadek zdegenerowany, wartości brzegowe. Testy własności: niezmienniczość na skalowanie, monotoniczność, odporność na permutację etykiet. Testy odporności: zerowa wariancja, współliniowość, wartości poza dziedziną, macierz osobliwa --- oczekiwanym wynikiem jest kontrolowany błąd z czytelnym komunikatem, a nie poprawny wynik. Podaj rezultat sprawdzenia zgodności na trzech systemach operacyjnych.

\wyroznienie{Narzędzia zewnętrzne.} Dla każdej biblioteki nazwa, źródło, przeznaczenie oraz krótki opis wejścia i wyjścia. Wymóg wprost ze Standardów.

\subsection{Jak pisać o pracy z narzędziem programistycznym}

Jeżeli korzystałeś z narzędzia wspomagającego pisanie kodu, opisz to jako element metodyki inżynierskiej, a nie jako wyznanie. Interesujące jest to, \wyroznienie{gdzie specyfikacja wychwyciła błąd}: w którym miejscu narzędzie uprościło założenie metody i jak to wykryłeś. To pokazuje zrozumienie metody lepiej niż bezbłędny kod. Zasada podstawowa pozostaje jedna: musisz umieć objaśnić każdą linię kodu, którą oddajesz.

\begin{tabelaepi}{Błędy w rozdziale inżynierskim}{tab:bledyimplementacja}
\begin{tabular}{@{}p{6.6cm}p{6.6cm}@{}}
\toprule
Błąd & Dlaczego szkodzi \\
\midrule
Wklejone całe pliki źródłowe & praca nie jest wydrukiem repozytorium \\
Opis funkcji zamiast opisu architektury & dokumentacja funkcji jest na stronie projektu \\
Procedura generowania danych zbyta jednym zdaniem & traci sens cała weryfikacja \\
Brak opisu kontraktu błędów & nie wiadomo, kiedy narzędzie zawodzi świadomie \\
Testy opisane jako napisano testy jednostkowe & nie wiadomo, co właściwie sprawdzono \\
\bottomrule
\end{tabular}
\end{tabelaepi}

\begin{kontrolna}
\item Schemat lub tabela modułów z odpowiedzialnościami
\item Opisany kontrakt: warunki wstępne, komunikaty błędów
\item Fragmenty kodu, nie całe pliki
\item Procedura generowania danych opisana z parametrami i ziarnem
\item Trzy warstwy weryfikacji opisane osobno
\item Wynik sprawdzenia zgodności na trzech systemach
\item Biblioteki zewnętrzne z nazwą, źródłem, przeznaczeniem, wejściem i wyjściem
\item Listingi zawierają wyłącznie znaki ASCII
\end{kontrolna}

\section{Rozdział trzeci: studium przypadku}

Plik \plik{rozdzialy/03-studium-przypadku.tex}, objętość osiem do dziesięciu stron.

Pokazujesz, że narzędzie działa na danych, a nie tylko się kompiluje. Zagadnienie ma należeć do dziedziny wskazanej w tytule pracy --- inaczej tytuł przestaje odpowiadać treści, co jest pierwszym kryterium oceny w Standardach.

\subsection{Co musi zawierać}

\wyroznienie{Charakterystyka danych.} Pochodzenie, licencja, liczba obserwacji i zmiennych, skale pomiarowe, braki. Jeżeli używasz danych z własnej procedury, odeślij do rozdziału drugiego i podaj przyjęte parametry. Napisz, jak wybór zbioru wpływa na to, co wolno wywnioskować.

Dane wygenerowane pozwalają sprawdzić poprawność, bo znasz prawdziwą odpowiedź. Dane empiryczne pozwalają pokazać użyteczność, ale prawdziwej odpowiedzi nie znasz. Najmocniejszy rozdział trzeci używa obu i mówi wprost, po co każdego.

\wyroznienie{Przebieg analizy.} Pełna ścieżka od danych surowych do wyniku, w postaci listingu. Czytelnik ma móc ją powtórzyć.

\wyroznienie{Wyniki.} Tabela i rysunek, a po nich interpretacja w tekście. Tabela nie zastępuje interpretacji. Oddziel to, co pokazuje narzędzie, od tego, co twierdzisz Ty: narzędzie zwraca liczby, a wniosek merytoryczny jest Twój i wymaga uzasadnienia.

\wyroznienie{Porównanie.} Zestaw wyniki z rezultatem z artykułu źródłowego albo z metodą odniesienia. Rozbieżność nie jest porażką --- jest wynikiem, o ile potrafisz wskazać jej przyczynę. Rozbieżność przemilczana jest błędem.

\begin{tabelaepi}{Błędy w studium przypadku}{tab:bledystudium}
\begin{tabular}{@{}p{6.6cm}p{6.6cm}@{}}
\toprule
Błąd & Dlaczego szkodzi \\
\midrule
Zbiór danych spoza dziedziny z tytułu pracy & treść przestaje odpowiadać tematowi \\
Tabela wyników bez interpretacji & czytelnik ma wykonać pracę autora \\
Wniosek mocniejszy niż dane & najczęstszy zarzut w recenzjach \\
Brak informacji o licencji danych & uniemożliwia odtworzenie analizy \\
Pominięcie rozbieżności z artykułem źródłowym & podważa wiarygodność weryfikacji \\
\bottomrule
\end{tabular}
\end{tabelaepi}

\begin{kontrolna}
\item Zbiór opisany: pochodzenie, licencja, rozmiar, skale, braki
\item Uzasadniony wybór danych względem tematu pracy
\item Pełna ścieżka analizy możliwa do powtórzenia
\item Każdy wynik zinterpretowany w tekście
\item Rozdzielone: co pokazuje narzędzie, a co twierdzi autor
\item Porównanie z wynikiem odniesienia wraz z omówieniem rozbieżności
\end{kontrolna}

\section{Podsumowanie}

Plik \plik{rozdzialy/04-podsumowanie.tex}, objętość dwie do trzech stron, bez numeru rozdziału.

Podsumowanie domyka klamrę z Wprowadzeniem. Czytelnik, który przeczyta tylko te dwie części, ma wiedzieć, co zostało zrobione i z jakim skutkiem. Zawiera: realizację celu podaną wprost, odpowiedź na pytanie badawcze wynikającą z rozdziałów drugiego i trzeciego, trzy do pięciu wniosków uporządkowanych od najważniejszego, ograniczenia oraz perspektywy rozwoju aplikacji. Ostatni element jest wymogiem wprost ze Standardów --- nie zbywaj go jednym zdaniem.

\wyroznienie{Nie wprowadzaj tu niczego nowego} i nie powołuj się na nowe źródła. Jeżeli w trakcie pisania podsumowania pojawia się nowy argument, jego miejsce jest wcześniej.

\begin{kontrolna}
\item Powrót do celu z Wprowadzenia
\item Odpowiedź na pytanie badawcze
\item Wnioski uporządkowane od najważniejszego
\item Ograniczenia
\item Perspektywy rozwoju opisane konkretnie
\item Zero nowych źródeł i nowych tez
\end{kontrolna}

\section{Elementy końcowe}

\subsection{Wykaz źródeł}\label{sec:wykazzrodel}

Plik \plik{rozdzialy/05-wykaz-zrodel.tex}. To \wyroznienie{nie to samo co bibliografia}. Bibliografia obejmuje opracowania naukowe, do których odsyłają powołania w tekście. Wykaz źródeł obejmuje narzędzia i dane: pakiety R, zbiory danych, serwisy udostępniające oprogramowanie. Standardy dopuszczają reprezentowanie narzędzi programistycznych adresami serwisów, które je udostępniają i opisują.

\subsection{Bibliografia}

Generowana automatycznie z pliku bibliograficznego. Zasada, której pilnuje recenzent: żadnej pozycji w bibliografii bez powołania w tekście i odwrotnie. Pozycje przeczytane, ale niewykorzystane, do bibliografii nie trafiają. Muszą się w niej znaleźć pozycje obcojęzyczne --- wymóg zarówno Standardów, jak i Instrukcji. Zapis wpisów opisuje rozdział~\ref{ch:bibliografia}.

\subsection{Spis ilustracji i indeks nazwisk}

Powstają automatycznie. Indeks zbiera nazwiska wstawione poleceniem \kod{\textbackslash osoba} --- używaj go od pierwszego dnia pisania, bo uzupełnianie indeksu na końcu to praca na kilka godzin.

\subsection{Aneksy}

Plik \plik{rozdzialy/06-aneksy.tex}. Wzór przewiduje cztery aneksy: metadane oprogramowania, notę wymaganą w serwisie internetowym zgodnie z załącznikiem do Standardów, oświadczenie o wykorzystaniu sztucznej inteligencji oraz wykaz poleceń wydanych narzędziu programistycznemu. Ostatni dotyczy wyłącznie budowy aplikacji i dla każdego polecenia podaje cel, treść, wprowadzone przez Ciebie zmiany i sposób weryfikacji. Aneksy muszą być powiązane z tekstem odsyłaczami.

\section{W jakiej kolejności pisać}\label{sec:kolejnosc}

\begin{tabelaepi}{Kolejność pisania rozdziałów}{tab:kolejnosc}
\begin{tabular}{@{}clp{7cm}@{}}
\toprule
Etap & Co piszesz & Dlaczego wtedy \\
\midrule
1 & rozdział 1, aparat formalny & zmusza do zrozumienia metody, zanim powstanie kod \\
2 & rozdział 1, założenia i granice & staje się kontraktem specyfikacji pakietu \\
3 & rozdział 2, generowanie danych & piszesz razem z kodem, gdy pamiętasz parametry \\
4 & rozdział 2, reszta & po zamknięciu implementacji \\
5 & rozdział 3 & po pierwszym pełnym przebiegu analizy \\
6 & rozdział 1, rozwiązania pokrewne & dopiero teraz wiesz, z czym się porównujesz \\
7 & Wprowadzenie & wiadomo już, co się zapowiada \\
8 & Podsumowanie & domyka klamrę z gotowym Wprowadzeniem \\
9 & aneksy, wykaz źródeł, redakcja & na końcu \\
\bottomrule
\end{tabular}
\end{tabelaepi}

Wprowadzenie pisane jako pierwsze jest prawie zawsze pisane dwa razy.

Trzy rzeczy prowadź \wyroznienie{od pierwszego dnia}, a nie na końcu: plik bibliograficzny, polecenia wstawiające nazwiska do indeksu oraz rejestr poleceń wydanych narzędziu programistycznemu. Odtwarzanie każdej z nich po fakcie kosztuje wielokrotnie więcej czasu.

\chapter{Warsztat pisania akademickiego}\label{ch:warsztat}

Rozdział jest o tym, \wyroznienie{jak} pisać: skąd wziąć pytanie, jak zbudować argument, jak rozmawiać z literaturą i jak nie stracić panowania nad źródłami. Co ma się znaleźć w której części pracy, opisuje rozdział~\ref{ch:jakpisac}; czego wymaga Instytut -- rozdział~\ref{ch:wymogi}.

\section{Pytanie analityczne}\label{sec:pytanie}

Celem pracy licencjackiej nie jest prezentacja wiedzy, lecz zrozumienie problemu. Proces ten obejmuje analizę dowodów i założeń oraz przemyślenie logiki argumentów. Punktem wyjścia jest pytanie -- i nie każde pytanie się do tego nadaje.

\subsection{Jakie pytanie nadaje się na fundament pracy}

Dobre pytanie analityczne dotyczy realnego dylematu obecnego w źródłach, a nie wymyślonego na potrzeby pracy. Daje odpowiedź, która nie jest oczywista przed przeprowadzeniem analizy. Sugeruje odpowiedź na tyle złożoną, że wymaga pełnej dyskusji, a nie akapitu. I da się na nie odpowiedzieć dostępnymi źródłami oraz dostępnymi danymi.

\subsection{Skąd brać pytania: punkty napięcia}

Pytania rodzą się w miejscach, gdzie źródła nie są ze sobą zgodne albo gdzie coś zgrzyta. Warto myśleć o nich jak o momentach, w których trzeba się zatrzymać i zastanowić, zanim da się iść dalej. Szukaj:

\begin{itemize}
\item stanowiska, które wydaje Ci się nieprzekonujące -- zapytaj, czego brakuje albo jak dowody można przeczytać inaczej;
\item nieścisłości, luki lub dwuznaczności w dowodach -- zapytaj, jak to zmienia rozumienie problemu;
\item nieoczekiwanego wniosku, który nie do końca się zgadza -- zapytaj, jak autorzy do niego doszli;
\item kontrowersji, która wymaga rozstrzygnięcia -- zapytaj, jak można ją rozstrzygnąć;
\item problemu, który został zignorowany -- zapytaj, dlaczego, albo spróbuj go rozwiązać;
\item dowodu, który zasługuje na bliższe przyjrzenie się.
\end{itemize}

\subsection{Napięcia właściwe dla pracy o oprogramowaniu badawczym}

W Twoim temacie napięcia są też innego rodzaju -- między metodą a jej użyciem. Artykuł podaje algorytm, ale \wyroznienie{nie precyzuje}, co zrobić z brakami danych, wartościami brzegowymi albo przypadkiem zdegenerowanym. Metoda zakłada dane, których w badaniach społecznych zwykle \wyroznienie{nie ma} w tej postaci. Istniejąca implementacja daje \wyroznienie{inny wynik} niż opis w publikacji. Autorzy pokazują metodę na danych symulowanych i nie mówią, co dzieje się na danych rzeczywistych. Procedura wymaga decyzji parametrycznej, którą publikacja przemilcza, choć wynik od niej zależy.

Każde z tych napięć jest gotowym punktem wyjścia dla pytania badawczego -- i jednocześnie decyzją projektową w pakiecie.

\subsection{Jak formułować}

Pytania \enquote{jak} i \enquote{dlaczego} wymagają więcej analizy niż \enquote{kto}, \enquote{co}, \enquote{gdzie} i \enquote{kiedy}. Dobre pytanie podkreśla wzorzec i związek albo sprzeczność i dylemat. Wyznacza zakres argumentacji, pozwalając skupić się na części szerokiego tematu, i może dotyczyć implikacji oraz konsekwencji analizy.

Pytanie prowadzi przez cały proces badawczy i pisanie. Kształtuje strukturę pracy, terminologię i kierunek poszukiwań -- dlatego warto poświęcić mu pierwsze tygodnie seminarium, a nie wymyślać je w tydzień przed oddaniem.

\section{Od pytania do tezy}

Pytanie otwiera, teza domyka. Teza to zdanie oznajmujące, z którym da się nie zgodzić. Jeżeli nikt rozsądny nie mógłby zaprzeczyć Twojej tezie, to nie jest teza, tylko stwierdzenie faktu.

\begin{tabelaepi}{Od pytania do tezy}{tab:teza}
\begin{tabular}{@{}p{5.2cm}p{8cm}@{}}
\toprule
Pytanie & Teza \\
\midrule
Jakie warunki musi spełnić implementacja metody, żeby wynik dał się odtworzyć? & Odtwarzalność wyniku zależy w większym stopniu od jawnego opisu obsługi braków danych niż od precyzji samego algorytmu. \\
Które założenia metody są najbardziej wrażliwe na jakość danych społecznych? & Metoda traci stabilność przy współliniowości kryteriów wcześniej niż przy brakach danych, co odwraca kolejność zaleceń podawanych w literaturze. \\
\bottomrule
\end{tabular}
\end{tabelaepi}

Teza może się zmienić w trakcie pracy. To normalne i nie jest porażką -- zmiana tezy pod wpływem wyników jest dowodem uczciwości, o ile praca opisuje wersję finalną, a nie obie naraz.

\section{Wejście w dyskusję akademicką}

Praca ma być głosem w rozmowie, która toczy się bez Ciebie. Są trzy sposoby włączenia się.

\subsection{Krok w przód}

Prezentujesz własną wiedzę i wkład: co zrobiłeś, czego to dowodzi, dlaczego to ważne. Wykazujesz znajomość literatury i podkreślasz, na czym polega unikalność Twojego rozwiązania. To dominujący tryb Wprowadzenia i rozdziału drugiego.

Zacznij od przeglądu istniejących prac, teorii i terminów. Czytaj nie tylko to, co autorzy mówią, ale też \wyroznienie{to, co pomijają albo traktują pobieżnie} -- tam zwykle leży luka. Następnie wskaż, gdzie dokładnie Twoja praca wpisuje się w te dyskusje. Nie wystarczy powiedzieć, że praca jest z tym związana; trzeba określić, jak Twoje podejście różni się od obecnego rozumienia albo je rozszerza.

\subsection{Krok w tył}

Umieszczasz własną argumentację w szerszym kontekście: jak praca koreluje z innymi badaniami, jak rozszerza albo kwestionuje istniejące ustalenia. To dominujący tryb rozdziału trzeciego i Podsumowania.

Zastanów się też, kto będzie Twoją pracą zainteresowany -- inni badacze, praktycy, szersza publiczność. Odbiorca decyduje o języku i o tym, które konsekwencje warto wyeksponować. Praca o oprogramowaniu badawczym ma odbiorcę podwójnego: metodologa, którego interesuje wierność implementacji, i badacza, którego interesuje, czy narzędzie da się użyć.

\subsection{Kwestionowanie}

Podważasz przyjęte założenia i pokazujesz alternatywne perspektywy. Wyjaśniasz, dlaczego te alternatywy są poznawczo ważne, ale niemożliwe do wykorzystania przy Twoim celu, i że stanowią obszar ograniczenia Twojej perspektywy, wykraczający poza Twój problem.

Bądź gotowy na kontrargumenty. To normalne, że recenzent nie zgodzi się z częścią wniosków. Przewidywanie sprzeciwów i odpowiadanie na nie w tekście wzmacnia pracę bardziej niż ich pominięcie.

\section{Zapożyczanie i rozwijanie teorii}

Zidentyfikuj teorie i pojęcia istotne dla tematu. Zwróć uwagę na kluczowe idee, definicje i argumenty -- bez tego nie da się ich skutecznie wprowadzić do własnej pracy.

W pracach z zakresu nauki o informacji pojawiają się terminy wymagające definicji, jak metadane, analiza semantyczna, uczenie maszynowe czy odtwarzalność. Wyjaśnij, jak używasz ich \wyroznienie{w kontekście swojej pracy} i jakie mają znaczenie dla zrozumienia problemu. Definicja robocza jest lepsza niż definicja słownikowa, o ile powiesz, skąd ją bierzesz.

Następnie rozwiń te idee: jak teorie stosują się do Twojego tematu, co trzeba w nich dostosować. Wreszcie pokaż, w jaki sposób te założenia wpłynęły na Twoje decyzje -- na kształt metod, na sposób weryfikacji, na interpretację wyników. Nie chodzi o powierzchowne odniesienia, lecz o widoczne konsekwencje.

W pracy projektowej ten mechanizm ma konkretną postać: \wyroznienie{założenie teoretyczne staje się warunkiem wstępnym funkcji}. Jeżeli metoda zakłada nieujemność ocen, w kodzie pojawia się sprawdzenie i komunikat błędu. Rozdział pierwszy i rozdział drugi mają się w tym miejscu spotkać -- i recenzent to sprawdzi.

\section{Czytanie artykułu metodycznego}

Artykuł, na którym opierasz pracę, czytasz inaczej niż literaturę przeglądową.

\wyroznienie{Kolejność.} Najpierw abstrakt i wnioski, żeby wiedzieć, po co to powstało. Potem sekcja definicyjna i algorytmiczna -- to jest Twój materiał. Sekcje dowodowe i porównawcze na końcu, i tylko w takim zakresie, w jakim wpływają na implementację.

\wyroznienie{Co wynotować.} Dla każdego wzoru: co jest wejściem, co wyjściem, jakiego typu i z jakiej dziedziny jest każdy symbol. Dla każdego kroku: co się dzieje, gdy wejście jest brzegowe. Osobno wypisz zdania zaczynające się od słów \enquote{zakładamy}, \enquote{przy założeniu}, \enquote{o ile} -- to jest przyszły kontrakt Twojego pakietu.

\wyroznienie{Czego szukać między wierszami.} Decyzji, które autorzy podjęli, ale ich nie uzasadnili: wartości domyślnych parametrów, sposobu normalizacji, kolejności kroków. Każda taka decyzja jest miejscem, w którym Twoja implementacja musi zająć stanowisko -- i w którym narzędzie wspomagające pisanie kodu najczęściej zgaduje.

\section{Notatki i zarządzanie źródłami}

Bibliografię prowadź \wyroznienie{od pierwszego dnia}. Odtwarzanie jej po fakcie zajmuje kilka dni i zawsze kończy się brakującym numerem strony.

Praktyka minimalna obejmuje cztery kroki. Utwórz kolekcję w Zotero z podkolekcjami na metodę, dziedzinę i narzędzia. Ustaw automatyczny eksport do pliku bibliograficznego projektu. Przy każdej pozycji zapisz jedno lub dwa zdania o tym, \wyroznienie{do czego ta pozycja jest Ci potrzebna w pracy} -- nie streszczenie, bo streszczenie jest w abstrakcie. Powołanie w tekście wstawiaj od razu przy pisaniu akapitu, a nie na końcu rozdziału.

Sprawdzanie metadanych po imporcie jest obowiązkowe: bazy bibliograficzne notorycznie gubią drugie imiona, mylą typ dokumentu i skracają tytuły czasopism. Szczegóły techniczne opisuje rozdział~\ref{ch:bibliografia}.

\section{Styl naukowy}

\subsection{Terminologia}

Termin raz wprowadzony ma jedną formę w całej pracy. Jeżeli w rozdziale pierwszym piszesz \enquote{wariant}, nie przechodź w rozdziale trzecim na \enquote{alternatywę}. Konsekwencja terminologiczna jest osobnym kryterium oceny \parencite{standardyEPI}.

Termin obcojęzyczny wprowadzasz raz, z polskim odpowiednikiem i zapisem oryginalnym w nawiasie, dalej używasz jednej formy. Nazw funkcji i pakietów nie odmieniamy -- poprzedzamy je rzeczownikiem: funkcja \fun{przygotuj\_dane}, pakiet \pkg{ggplot2}.

\subsection{Akapit}

Akapit ma jedną myśl i mieści się na jednej trzeciej strony. Pierwsze zdanie mówi, o czym akapit będzie. Ostatnie wiąże go z następnym. Akapit na całą stronę zawsze da się podzielić i zawsze na tym zyskuje.

\subsection{Rejestr}

Praca naukowa różni się od publicystyki i od dokumentacji technicznej.

\begin{tabelaepi}{Rejestr wypowiedzi}{tab:rejestr}
\begin{tabular}{@{}p{6.4cm}p{6.8cm}@{}}
\toprule
Zamiast & Napisz \\
\midrule
Niesamowicie ważnym zagadnieniem jest\ldots & Zagadnienie ma znaczenie dla\ldots, ponieważ\ldots \\
Jak wiadomo, media społecznościowe\ldots & Badania nad mediami społecznościowymi wskazują, że\ldots (powołanie) \\
Postanowiłem użyć tej metody, bo jest wygodna & Wybrano metodę ze względu na\ldots \\
Funkcja bierze dane i je przetwarza & Funkcja przyjmuje macierz ocen i zwraca wektor wag \\
\bottomrule
\end{tabular}
\end{tabelaepi}

Unikaj nadmiernych skrótów myślowych, wypowiedzi emocjonalnych i publicystycznych, rozbudowanych dygresji oraz treści luźno związanych z tokiem rozumowania. Nie opieraj wniosków na nieujawnionych przesłankach ani na niepotwierdzonych naukowo przekonaniach \parencite{instrukcjaISI}.

\subsection{Precyzja}

Tezy, sądy i wnioski formułuj precyzyjnie. Zdanie o tym, że metoda działa lepiej, nie znaczy nic; zdanie o tym, że metoda daje niższy błąd średniokwadratowy przy próbach poniżej stu obserwacji, znaczy. Każde twierdzenie porównawcze wymaga wskazania, względem czego i na jakiej podstawie.

\section{Rzetelność}

\subsection{Oryginalność}

Praca nie może powielać badań przeprowadzonych wcześniej przez inne osoby -- z wyjątkiem sytuacji, gdy chodzi o ich weryfikację lub aktualizację. Implementacja opublikowanej metody \wyroznienie{jest} dopuszczalna i mieści się w tym wyjątku, o ile praca wnosi własny wkład: rozstrzygnięcie niedopowiedzeń, procedurę weryfikacji, zastosowanie w nowym obszarze.

\subsection{Kompilacja}

Praca ani żaden jej fragment nie może być zestawieniem zapożyczeń -- ani dosłownych cytatów, ani fragmentów sparafrazowanych. Zasada obejmuje także tłumaczenia z języka obcego. Parafraza nie zwalnia z powołania: wyrażenie cudzej myśli własnymi słowami bez wskazania autora jest naruszeniem, nawet jeżeli żadne zdanie nie zostało przepisane.

Dosłowne cytowanie i referowanie cudzych ustaleń ogranicz do minimum koniecznego do ustalenia stanu badań, podjęcia z nimi dyskusji albo osadzenia własnych wyników.

\subsection{Granica między własnym a cudzym}

Wypowiedź musi być tak skonstruowana, a przypisy umieszczone w takich miejscach, żeby jednoznacznie oddzielić własne poglądy od zapożyczonych. \wyroznienie{W żadnym miejscu czytelnik nie może mieć wątpliwości, co jest wynikiem Twojej pracy, a co odwołaniem do cudzej.} W praktyce oznacza to zwroty sygnalizujące: cudze wprowadzasz zwrotami w rodzaju \enquote{zdaniem}, \enquote{jak wykazuje}, \enquote{w ujęciu}; własne -- zwrotami \enquote{w niniejszej pracy przyjęto}, \enquote{uzyskane wyniki wskazują}, \enquote{na tej podstawie można stwierdzić}.

Zasada nie dotyczy wiedzy powszechnej, danych powszechnie dostępnych i ogólnie określonych metod badawczych. Dotyczy natomiast konkretnych teorii, idei, danych i narzędzi badawczych, które stanowią dorobek konkretnych osób.

\subsection{Stan badań w układzie problemowym}

Dotychczasowy dorobek poddajesz samodzielnej analizie i przedstawiasz w postaci własnych wniosków, uporządkowanych \wyroznienie{problemowo, a nie chronologicznie ani alfabetycznie}. Układ docelowy: co już wiemy i co jeszcze wymaga zbadania.

Przegląd, który brzmi jak wyliczanka kolejnych autorów i tego, co napisali, nie jest analizą, tylko listą. Analiza brzmi inaczej: w danej kwestii literatura jest zgodna co do jednego, natomiast rozchodzi się w sprawie drugiego, gdzie jedni autorzy wskazują na to, a inni przeciwnie.

\section{Redakcja i korekta}

Redakcja to osobny etap, a nie przedłużenie pisania. Zaplanuj na nią co najmniej tydzień.

Przeglądy rób po kolei, za każdym razem szukając jednej rzeczy. Najpierw struktura: czy kolejność rozdziałów i sekcji odpowiada logice wywodu. Potem argument: czy każda teza ma uzasadnienie, a każde uzasadnienie prowadzi do tezy. Następnie powołania: czy każde twierdzenie o cudzym dorobku ma przypis i czy każda pozycja bibliografii ma powołanie. Dalej terminologia: czy nazwy są konsekwentne w całej pracy. Potem język: zdania, akapity, interpunkcja. Na końcu skład: wiszące wiersze, pływające tabele, przeniesienia, spisy.

Korektę rób na wydruku albo w innym kroju pisma. Czytaj na głos zdania, które wydają się długie; zdanie, którego nie da się przeczytać jednym tchem, wymaga podziału. Przed ostatnim przeglądem odłóż tekst na dwa lub trzy dni -- nie ma sposobu, żeby to obejść.

\section{Harmonogram}

Pisanie i budowa pakietu idą równolegle. Rozdziały powstają w kolejności innej niż kolejność w pracy; podaje ją tabela~\ref{tab:kolejnosc}.

\begin{tabelaepi}{Co prowadzić od pierwszego dnia}{tab:odpoczatku}
\begin{tabular}{@{}p{4.2cm}p{3.8cm}p{5.6cm}@{}}
\toprule
Co & Gdzie & Dlaczego \\
\midrule
plik bibliograficzny & katalog bibliografii & odtwarzanie po fakcie zajmuje kilka dni \\
polecenia wstawiające nazwiska & rozdziały & indeks nazwisk powstaje sam \\
rejestr poleceń dla narzędzia & repozytorium pakietu & wymagany w aneksie, nie da się odtworzyć z pamięci \\
\bottomrule
\end{tabular}
\end{tabelaepi}

Reguła praktyczna na koniec: \wyroznienie{nie pisz akapitu, którego nie umiesz opowiedzieć na głos w trzech zdaniach}. Jeżeli nie umiesz, nie jest to problem pisania, tylko znak, że ta część nie jest jeszcze przemyślana.

\chapter{Wymogi formalne}\label{ch:wymogi}

Zestawienie wymagań obowiązujących pracę licencjacką na kierunku elektroniczne przetwarzanie informacji, z podaniem źródła każdego z nich. Dwa dokumenty Instytutu uzupełniają się i miejscami rozchodzą -- przy rozbieżności o \wyroznienie{układzie pracy} rozstrzygają Standardy \parencite{standardyEPI}, a o \wyroznienie{zapisie bibliografii} Instrukcja \parencite{instrukcjaISI}. W tabelach skrót S oznacza Standardy, a skrót I -- Instrukcję.

\section{Skład}

\begin{tabelaepi}{Wymogi składu}{tab:sklad}
\begin{tabular}{@{}p{8.4cm}cp{3.4cm}@{}}
\toprule
Wymóg & Źródło & Kto realizuje \\
\midrule
Pismo 12 pkt w tekście głównym & S & klasa \\
Interlinia 1,5 & S & klasa \\
Marginesy 2,5 cm z każdej strony & S, I & klasa \\
Przy oprawie margines wewnętrzny szerszy kosztem zewnętrznego, nie mniej niż 1,5 cm & I & opcja klasy \\
Tytuły rozdziałów tym samym krojem, większe i wytłuszczone & S & klasa \\
Tytuły rozdziałów 18 pkt, wyśrodkowane & I & klasa \\
Tytuły podrozdziałów 14 pkt, wytłuszczone, do lewej & I & klasa \\
Numeracja rozdziałów wielopoziomowa & S, I & klasa \\
Wstęp i Zakończenie bez numeracji, ale w spisie treści & I & klasa \\
Każdy element ze spisu treści od nowej strony & S, I & klasa \\
Wszystkie strony numerowane & S & klasa \\
Numer strony w prawym dolnym rogu & I & klasa \\
Strona tytułowa liczona, bez numeru & S, I & klasa \\
Tekst wyjustowany & S, I & klasa \\
\bottomrule
\end{tabular}
\end{tabelaepi}

Wszystkie pozycje realizuje klasa dokumentu bez Twojego udziału. Nie zmieniaj ustawień składu w preambule -- rozbieżność z wymogami obciąża pracę, a nie klasę.

\section{Objętość}

Tekst od pierwszej strony Wprowadzenia do ostatniej strony Podsumowania nie może przekraczać \wyroznienie{czterdziestu stron znormalizowanych, czyli 72 000 znaków ze spacjami} \parencite{instrukcjaISI}. Do limitu nie wliczają się: strona tytułowa, strony informacyjne, spis treści, wykaz źródeł, bibliografia, spisy, indeks i aneksy.

Kontrolę wykonuje skrypt \plik{tools/policz\_znaki.R}, uruchamiany z katalogu wzoru. Liczy znaki po usunięciu poleceń, listingów, wzorów i tabel, więc daje oszacowanie od góry.

\section{Struktura}

Kolejność elementów we wzorze łączy wymagania obu dokumentów: strona tytułowa, opcjonalna strona informacyjna z abstraktem i słowami kluczowymi, spis treści, Wprowadzenie, rozdziały numerowane, Podsumowanie, wykaz źródeł, bibliografia, spis ilustracji, indeks nazwisk, aneksy.

Standardy opisują treść rozdziałów zasadniczych jako logikę aplikacji -- schemat ogólny, moduły, opis interfejsu ze zrzutami ekranu -- oraz implementację, czyli format danych, strukturę bazy, sposób realizacji modułów, biblioteki zewnętrzne i narzędzia. Wzór realizuje to w trzech rozdziałach.

\subsection{Strona informacyjna}

Wymóg Instrukcji, włączany opcją klasy. Zawiera opis bibliograficzny pracy w postaci obejmującej nazwisko i imię autora, rok obrony, tytuł, oznaczenie rodzaju pracy, promotora, miejsce, jednostkę i liczbę stron. Dalej abstrakt liczący \wyroznienie{od 1000 do 1500 znaków ze spacjami}, przedstawiający cel, przedmiot, metody i kluczowe wyniki, napisany językiem skondensowanym i bez ogólników. Na końcu \wyroznienie{około pięciu słów kluczowych} w porządku alfabetycznym, rozdzielonych przecinkami.

Strona angielska powstaje automatycznie po wypełnieniu angielskich metadanych i zawiera dokładnie te same informacje w tym samym układzie.

\section{Cytowanie}

Standardy wymieniają pięć sytuacji, w których wskazanie źródła jest obowiązkowe:

\begin{enumerate}
\item cytat bezpośredni, czyli dosłownie przytoczone cudze słowa;
\item graficzna forma cytatu -- cudze tabele, rysunki, zestawienia;
\item cytat znany pośrednio, z tekstu jeszcze innego autora;
\item zebrane przez kogoś informacje: ankiety, dane statystyczne, dokumenty, materiały ze stron internetowych;
\item cudza teza, argument, opinia, idea, interpretacja, unikalne ujęcie tematu.
\end{enumerate}

Każde przytoczenie umieszczasz w cudzysłowie polskim, a cytat w cytacie w cudzysłowie wewnętrznym. Wewnątrz cytatu obowiązuje pisownia cytowanego autora; zmianę odnotowujesz w nawiasie kwadratowym po cudzysłowie zamykającym. Tłumaczenie własne oznaczasz w nawiasie kwadratowym, a cytat w oryginale przywołujesz w przypisie dolnym. Fragment powyżej dwóch lub trzech zdań wyodrębniasz graficznie.

Przytoczenie musi być \wyroznienie{wkomponowane w tekst}: wprowadzone lub domknięte komentarzem, który sygnalizuje obecność cudzego słowa. Samo wklejenie z przypisem nie wystarcza. Przy powoływaniu się na cudzą opinię obowiązuje, niezależnie od przypisu, wprowadzenie z informacją, do kogo dany pogląd należy.

Brak zasygnalizowania w tekście, że posługujesz się cudzym słowem, oznacza naruszenie prawa \parencite{prawoAutorskie}. Brak odpowiednio opisanych źródeł może być powodem dyskwalifikacji pracy.

\section{Przypisy i bibliografia}

Obowiązuje system nawiasowy, tak zwany harwardzki: nazwisko, rok, w razie potrzeby strona. Postacie powołań i odpowiadające im polecenia zestawia tabela~\ref{tab:powolania}. Skracanie listy autorów, dodawanie inicjału i sufiksów rocznika wykonuje styl automatycznie -- Twoim zadaniem jest poprawny wpis w pliku bibliograficznym.

Do przypisów dolnych trafiają wyłącznie przypisy dygresyjne, których należy unikać, oraz oryginalne brzmienie tłumaczonych cytatów.

Zasada wiążąca bibliografię z tekstem: \wyroznienie{bibliografia zawiera wyłącznie pozycje faktycznie wykorzystane, do których odsyłają przypisy w tekście}. Pozycje dotyczące tematu, ale niewykorzystane, do bibliografii nie trafiają. Muszą się w niej znaleźć pozycje obcojęzyczne.

\section{Materiał ilustracyjny}

\begin{tabelaepi}{Wymogi wobec materiału ilustracyjnego}{tab:ilustracje}
\begin{tabular}{@{}p{10.6cm}c@{}}
\toprule
Wymóg & Źródło \\
\midrule
Numeracja ciągła przez cały tekst, osobna dla każdego typu & I \\
Tytuł nad ilustracją & I \\
Źródło pod ilustracją & I \\
Oznaczenie opracowania własnego, także na podstawie cudzego materiału & I \\
Pełny opis bibliograficzny dla materiału przejętego & I \\
Interfejs aplikacji zilustrowany zrzutami ekranu & S \\
\bottomrule
\end{tabular}
\end{tabelaepi}

Środowiska wzoru realizują tytuł, numerację i miejsce na źródło automatycznie.

Prawa do ilustracji: bez zgody autora można wykorzystać materiały z domeny publicznej, na licencji Creative Commons lub podobnej oraz zrzuty ekranu z publicznie dostępnych źródeł elektronicznych. Pozostałe wymagają pisemnej zgody właściciela autorskich praw majątkowych.

\section{Zapis w tekście}

\begin{tabelaepi}{Zapis elementów w tekście}{tab:zapiselementow}
\begin{tabular}{@{}p{6.8cm}ll@{}}
\toprule
Element & Zapis & Polecenie \\
\midrule
tytuł czasopisma, książki, zasobu sieciowego & kursywa bez cudzysłowu & \kod{\textbackslash termin} \\
wyraz obcojęzyczny & kursywa & \kod{\textbackslash termin} \\
wyróżnienie treściowe & wytłuszczenie & \kod{\textbackslash wyroznienie} \\
cytat & cudzysłów polski & \kod{\textbackslash enquote} \\
nazwa funkcji & pismo maszynowe & \kod{\textbackslash fun} \\
nazwa argumentu & pismo maszynowe & \kod{\textbackslash argument} \\
nazwa pakietu & pismo maszynowe & \kod{\textbackslash pkg} \\
nazwa pliku & pismo maszynowe & \kod{\textbackslash plik} \\
\bottomrule
\end{tabular}
\end{tabelaepi}

\section{Wymogi projektu dyplomowego}

Poza pracą pisemną Standardy wymagają samodzielnie zaprojektowanej i wykonanej aplikacji z interfejsem internetowym. Zaliczenie seminarium następuje dopiero po przyjęciu projektu \parencite{standardyEPI,prawoSzkolnictwo}.

Aplikacja nastawiona na prezentację informacji musi robić więcej niż wyświetlanie zawartości bazy danych: konieczny jest niestandardowy pomysł prezentacji, dostosowany projekt graficzny i mechanizm interakcji z użytkownikiem.

W serwisie obowiązkowa jest nota z załącznika do Standardów, umieszczona w zakładce o serwisie. Podaje ona autora serwisu, informację, że serwis stanowi integralną część pracy licencjackiej na kierunku elektroniczne przetwarzanie informacji, oraz promotora i jednostkę. Ta sama nota trafia do aneksu pracy.

\section{Kryteria oceny}

Standardy wymieniają wprost, co podlega ocenie promotora i recenzenta: zgodność treści pracy z tematem, strukturę pracy, znajomość stanu badań lub praktyki w zakresie wyznaczonym przez tematykę, poprawność językową i terminologiczną, metodologię projektu, zawartość merytoryczną, poprawność działania oraz estetykę aplikacji, a także dobór i wykorzystanie piśmiennictwa. Recenzje trafiają do Archiwum Prac.

\section{Lista kontrolna przed oddaniem}\label{sec:listakontrolna}

\subsection*{Kompilacja i objętość}

\begin{kontrolna}
\item Pełna kompilacja kończy się bez błędów
\item Brak komunikatów o nierozwiązanych powołaniach i odsyłaczach
\item Objętość mieści się w limicie
\item Listingi nie zawierają znaków spoza ASCII
\end{kontrolna}

\subsection*{Treść}

\begin{kontrolna}
\item Tytuł pracy odpowiada zawartości i zawiera zagadnienie ogólniejsze
\item Cel z Wprowadzenia domknięty w Podsumowaniu
\item Perspektywy rozwoju aplikacji opisane konkretnie
\item Wszystkie przytoczenia wprowadzone komentarzem w tekście
\item Własne ustalenia oddzielone od cudzych na poziomie zdania
\end{kontrolna}

\subsection*{Bibliografia i źródła}

\begin{kontrolna}
\item Każda pozycja bibliografii ma powołanie w tekście
\item Każde powołanie ma pozycję w bibliografii
\item W bibliografii są pozycje obcojęzyczne
\item Wykaz źródeł obejmuje wszystkie użyte pakiety
\item Adresy sieciowe mają datę odczytu
\end{kontrolna}

\subsection*{Elementy końcowe}

\begin{kontrolna}
\item Spis ilustracji zgodny z zawartością pracy
\item Indeks nazwisk obejmuje nazwiska z tekstu i bibliografii
\item Aneksy powiązane z tekstem odsyłaczami
\item Aneks z metadanymi oprogramowania wypełniony
\item Aneks z notą serwisu zgodny z treścią na stronie projektu
\item Oświadczenie o wykorzystaniu sztucznej inteligencji wypełnione, jeżeli dotyczy
\item Wykaz poleceń wydanych narzędziu programistycznemu kompletny
\end{kontrolna}

\subsection*{Projekt}

\begin{kontrolna}
\item Strona projektu działa pod adresem podanym na stronie tytułowej
\item Nota z załącznika obecna w serwisie
\item Repozytorium publiczne i zgodne z adresem ze strony tytułowej
\item Sprawdzenie pakietu przechodzi na trzech systemach operacyjnych
\end{kontrolna}

\subsection*{Redakcja}

\begin{kontrolna}
\item Korekta całości po wydruku, nie na ekranie
\item Tekst wyjustowany, bez wiszących wierszy
\item Konsekwentna terminologia w całej pracy
\item Wyszukanie uwag promotora w plikach źródłowych nie daje wyników
\item Usunięte własne komentarze robocze i wskazówki wzoru
\end{kontrolna}

\input{rozdzialy/07-agent-cli}
\input{rozdzialy/08-etyka-si}

\bibliografia
\spisilustracji

\end{document}
